\documentclass{article}

\usepackage[final]{neurips_2025}
\makeatletter\renewcommand{\@neuripsyear}{2026}\makeatother

\usepackage[utf8]{inputenc}
\usepackage[T1]{fontenc}
\usepackage{lmodern}
\usepackage{hyperref}
\usepackage{url}
\usepackage{booktabs}
\usepackage{amsfonts}
\usepackage{amsmath}
\usepackage{amssymb}
\usepackage{amsthm}
\usepackage{nicefrac}
\usepackage{microtype}
\usepackage{xcolor}
\usepackage{colortbl}
\usepackage{graphicx}
\usepackage{subcaption}
\usepackage{multirow}
\usepackage{algorithm}
\usepackage{algorithmic}
\usepackage{tablefootnote}
\usepackage{enumitem}
\usepackage{tikz}
\usetikzlibrary{shapes,arrows,positioning,fit,calc}

\definecolor{improve}{RGB}{0,128,0}
\definecolor{decline}{RGB}{200,0,0}
\definecolor{easycolor}{RGB}{76,175,80}
\definecolor{medcolor}{RGB}{255,152,0}
\definecolor{hardcolor}{RGB}{244,67,54}
\newcommand{\up}[1]{\textcolor{improve}{\textbf{+#1}}}
\newcommand{\down}[1]{\textcolor{decline}{\textbf{#1}}}
\newcommand{\ci}[2]{\scriptsize[#1,\,#2]}
\DeclareMathOperator*{\argmax}{arg\,max}
\DeclareMathOperator*{\argmin}{arg\,min}
\newcommand{\method}{\textsc{ReasonSpec}}
\newcommand{\fullmethod}{Reasoning Speculation}

\newtheorem{proposition}{Proposition}
\newtheorem{corollary}{Corollary}
\theoremstyle{remark}
\newtheorem{remark}{Remark}

% NeurIPS checklist
\providecommand{\answerYes}[1]{\textbf{Yes.} #1}
\providecommand{\answerNo}[1]{\textbf{No.} #1}
\providecommand{\answerNA}[1]{\textbf{N/A.} #1}

\title{\fullmethod{}: Parallel Exploration with Cross-Path\\Consensus for Adaptive Test-Time Compute}

\author{
  Anonymous Author(s)
}

\begin{document}

\maketitle

\begin{abstract}
% PLACEHOLDER — will be finalized after experiments complete
Scaling test-time compute via chain-of-thought reasoning improves accuracy but wastes tokens on easy problems and underallocates on hard ones.
Existing adaptive methods rely on single-path difficulty signals (answer presence, uncertainty, learned features) which we show are fundamentally unreliable: all such black-box controllers fail to outperform fixed budgets.
We introduce \method{}, a training-free framework that restructures the computation of reasoning itself.
Instead of allocating a single token budget and hoping it matches problem difficulty, \method{} generates $K$ short parallel exploration paths, computes a \emph{cross-path consensus signal}, and adaptively routes each problem:
high consensus $\rightarrow$ accept the majority answer (easy); partial consensus $\rightarrow$ extend the most promising path (medium); low consensus $\rightarrow$ full deliberation with all paths as context (hard).
The key insight is that multi-path agreement is a far more reliable difficulty signal than any feature extractable from a single reasoning trace (Spearman $\rho = -0.51$ vs.\ difficulty, $p < 10^{-88}$).
On GSM8K with Qwen3-8B, \method{} achieves \textbf{XX.X\%} accuracy using only \textbf{XXX} average tokens, compared to \textbf{XX.X\%} for a fixed budget using the same total compute, representing a \textbf{+X.X\,pp} improvement at \textbf{XX\%} token savings.
\method{} is model-agnostic, requires no training, and works at the API level with any reasoning LLM.
\end{abstract}

% ============================================================================
\section{Introduction}
\label{sec:intro}
% ============================================================================

% intro_reasonspec.tex — Introduction for the Reasoning Speculation paper
% Loaded by main_reasonspec.tex via % intro_reasonspec.tex — Introduction for the Reasoning Speculation paper
% Loaded by main_reasonspec.tex via % intro_reasonspec.tex — Introduction for the Reasoning Speculation paper
% Loaded by main_reasonspec.tex via \input{sections/intro_reasonspec}

Large language models equipped with chain-of-thought reasoning have achieved remarkable progress on complex tasks by ``thinking longer''---generating extended reasoning traces before committing to an answer.
This paradigm, exemplified by OpenAI o1~\cite{openai2024o1}, QwQ~\cite{qwq2024}, and DeepSeek-R1~\cite{deepseekr1}, frames test-time performance as a scaling problem: allocate more tokens, and accuracy improves.
Yet the gains from extended reasoning are far from uniform across problem instances.
Easy problems are \emph{overthought}: the model wastes hundreds of tokens re-deriving trivial facts or, worse, ``talks itself out of'' a correct intermediate answer through unnecessary self-revision.
Hard problems are \emph{underthought}: a fixed budget truncates reasoning before the model reaches the critical insight.
The fundamental mismatch between uniform compute allocation and heterogeneous problem difficulty leaves substantial accuracy and efficiency gains on the table.

\paragraph{The overthinking problem is pervasive.}
Through systematic experiments across GSM8K with Qwen3-8B, we quantify the extent of this mismatch.
When the reasoning budget is increased from a short to a long allocation, \textbf{XX.X\%} of samples exhibit \emph{accuracy degradation}---they were answered correctly at a lower budget but incorrectly at a higher one.
This is not an isolated edge case; it is a structured phenomenon affecting roughly one-third of all instances.
The ideal solution is clear in principle: allocate compute adaptively, spending little on easy problems and much on hard ones.
An oracle that selects the optimal budget per instance achieves \textbf{XX.X\%} accuracy---far above any fixed budget---demonstrating the enormous headroom for adaptive methods.

\paragraph{Existing adaptive approaches fail in the black-box setting.}
If the gains from adaptive allocation are so large, why has no practical method captured them?
We conducted what is, to our knowledge, the first comprehensive evaluation of black-box adaptive budget controllers for reasoning LLMs.
We tested all major categories of single-path difficulty signals:
(i)~\emph{honest features} extracted from a partial reasoning trace (token utilization, answer presence, lexical coherence), yielding controllers at $-6.5$\,pp below the fixed baseline;
(ii)~\emph{uncertainty-based routing} using answer entropy or self-reported confidence, at $-6.4$\,pp;
and (iii)~\emph{MCTS-style exploration--refinement}, at $-5.6$\,pp.
\textbf{All single-path controllers failed}, often performing \emph{worse} than the na\"ive fixed budget they were designed to improve.
The root cause is informational: a single reasoning trace is fundamentally impoverished as a difficulty signal.
The model's own confidence is poorly calibrated, token-level statistics are noisy, and partial traces reveal little about whether the \emph{final} answer will be correct.
In short, one witness is not enough to judge whether a problem is easy or hard.

\paragraph{Key insight: change the computation structure.}
Rather than trying to extract a weak difficulty signal from a single reasoning path, we propose to \emph{restructure the computation itself} to produce a fundamentally richer signal.
The core idea is simple: instead of running one long chain of thought, first run $K$ short, independent reasoning paths in parallel.
The \emph{cross-path consensus}---the degree to which these $K$ independent ``witnesses'' agree on an answer---provides a difficulty signal that is categorically stronger than anything extractable from a single trace.
Intuitively, if $K$ independent short explorations all converge to the same answer, the problem is almost certainly easy; if they scatter across many answers, it is hard and warrants extended deliberation.
This is analogous in spirit to how speculative decoding~\cite{leviathan2023fast} changed the \emph{structure} of the generation process (small model drafts, large model verifies) without changing the model itself---here, we change the structure of reasoning (parallel exploration, consensus-based routing) without changing the model or its training.
We show in \S\ref{sec:theory} that multi-path agreement provides an information-theoretically optimal difficulty estimator: when each path has independent error probability $\varepsilon$, the consensus signal concentrates exponentially in $K$, yielding a Spearman correlation with true difficulty of $\rho = -0.51$ ($p < 10^{-88}$)---far stronger than any single-path feature we tested.

\paragraph{The \method{} framework.}
We introduce \textbf{\fullmethod{}} (\method{}), a training-free, model-agnostic framework for adaptive test-time compute that operates in three phases:
\begin{enumerate}[nosep,leftmargin=*]
    \item \textbf{Explore.} Given a question $x$, generate $K$ short reasoning paths $\{r_1, \ldots, r_K\}$, each with a small token budget $b_{\text{explore}}$. These paths are independent and can be parallelized.
    \item \textbf{Decide.} Extract the answer from each path and compute a consensus score $c(x) = \max_a \frac{1}{K} \sum_{k=1}^{K} \mathbf{1}[a_k = a]$. Route the problem based on $c(x)$: high consensus $\rightarrow$ accept the majority answer immediately (easy); partial consensus $\rightarrow$ extend the most promising path with additional tokens (medium); low consensus $\rightarrow$ invoke full deliberation with all $K$ paths as context (hard).
    \item \textbf{Resolve.} For medium and hard problems, generate extended reasoning conditioned on the exploration paths, using a budget calibrated to the estimated difficulty tier.
\end{enumerate}
\method{} requires no training, no access to model internals, and no modification to the underlying LLM.
It operates entirely at the API level and is compatible with any reasoning model that supports temperature-based sampling.
The total compute budget is controlled by $K$, $b_{\text{explore}}$, and the per-tier continuation budgets, yielding a flexible accuracy--cost tradeoff.

\paragraph{Contributions.}
Our work makes the following contributions:
\begin{itemize}[nosep,leftmargin=*]
    \item \textbf{A new computation paradigm for adaptive test-time compute.}
    We introduce parallel exploration with consensus-based routing as an alternative to single-path budget allocation.
    Unlike prior adaptive methods that extract features from one reasoning trace, \method{} restructures the computation to produce a multi-path consensus signal that is fundamentally more informative about problem difficulty.

    \item \textbf{Information-theoretic justification.}
    We prove that multi-path consensus dominates any single-path difficulty feature in the black-box setting (\S\ref{sec:theory}).
    The consensus signal concentrates exponentially in $K$, explaining both why single-path controllers fail and why even modest $K$ (e.g., $K=5$) suffices.

    \item \textbf{Comprehensive empirical evaluation.}
    On GSM8K with Qwen3-8B, \method{} achieves \textbf{XX.X\%} accuracy at \textbf{XX\%} of the token cost of a fixed high-budget baseline, representing a \textbf{+X.X\,pp} improvement at matched compute and \textbf{XX\%} token savings at matched accuracy.
    We provide full ablations over $K$, $b_{\text{explore}}$, consensus thresholds, and continuation strategies.

    \item \textbf{First systematic negative results for black-box controllers.}
    We show that \emph{all} major categories of single-path adaptive controllers (honest features, uncertainty routing, MCTS refinement) fail to outperform fixed budgets, with degradations of $5$--$7$\,pp.
    These negative results are themselves a contribution, delimiting what is achievable without structural changes to the computation.

    \item \textbf{Training-free and model-agnostic.}
    \method{} requires no fine-tuning, no reward models, and no access to logits or hidden states.
    It works with any LLM API that supports standard generation and is immediately deployable.
\end{itemize}


Large language models equipped with chain-of-thought reasoning have achieved remarkable progress on complex tasks by ``thinking longer''---generating extended reasoning traces before committing to an answer.
This paradigm, exemplified by OpenAI o1~\cite{openai2024o1}, QwQ~\cite{qwq2024}, and DeepSeek-R1~\cite{deepseekr1}, frames test-time performance as a scaling problem: allocate more tokens, and accuracy improves.
Yet the gains from extended reasoning are far from uniform across problem instances.
Easy problems are \emph{overthought}: the model wastes hundreds of tokens re-deriving trivial facts or, worse, ``talks itself out of'' a correct intermediate answer through unnecessary self-revision.
Hard problems are \emph{underthought}: a fixed budget truncates reasoning before the model reaches the critical insight.
The fundamental mismatch between uniform compute allocation and heterogeneous problem difficulty leaves substantial accuracy and efficiency gains on the table.

\paragraph{The overthinking problem is pervasive.}
Through systematic experiments across GSM8K with Qwen3-8B, we quantify the extent of this mismatch.
When the reasoning budget is increased from a short to a long allocation, \textbf{XX.X\%} of samples exhibit \emph{accuracy degradation}---they were answered correctly at a lower budget but incorrectly at a higher one.
This is not an isolated edge case; it is a structured phenomenon affecting roughly one-third of all instances.
The ideal solution is clear in principle: allocate compute adaptively, spending little on easy problems and much on hard ones.
An oracle that selects the optimal budget per instance achieves \textbf{XX.X\%} accuracy---far above any fixed budget---demonstrating the enormous headroom for adaptive methods.

\paragraph{Existing adaptive approaches fail in the black-box setting.}
If the gains from adaptive allocation are so large, why has no practical method captured them?
We conducted what is, to our knowledge, the first comprehensive evaluation of black-box adaptive budget controllers for reasoning LLMs.
We tested all major categories of single-path difficulty signals:
(i)~\emph{honest features} extracted from a partial reasoning trace (token utilization, answer presence, lexical coherence), yielding controllers at $-6.5$\,pp below the fixed baseline;
(ii)~\emph{uncertainty-based routing} using answer entropy or self-reported confidence, at $-6.4$\,pp;
and (iii)~\emph{MCTS-style exploration--refinement}, at $-5.6$\,pp.
\textbf{All single-path controllers failed}, often performing \emph{worse} than the na\"ive fixed budget they were designed to improve.
The root cause is informational: a single reasoning trace is fundamentally impoverished as a difficulty signal.
The model's own confidence is poorly calibrated, token-level statistics are noisy, and partial traces reveal little about whether the \emph{final} answer will be correct.
In short, one witness is not enough to judge whether a problem is easy or hard.

\paragraph{Key insight: change the computation structure.}
Rather than trying to extract a weak difficulty signal from a single reasoning path, we propose to \emph{restructure the computation itself} to produce a fundamentally richer signal.
The core idea is simple: instead of running one long chain of thought, first run $K$ short, independent reasoning paths in parallel.
The \emph{cross-path consensus}---the degree to which these $K$ independent ``witnesses'' agree on an answer---provides a difficulty signal that is categorically stronger than anything extractable from a single trace.
Intuitively, if $K$ independent short explorations all converge to the same answer, the problem is almost certainly easy; if they scatter across many answers, it is hard and warrants extended deliberation.
This is analogous in spirit to how speculative decoding~\cite{leviathan2023fast} changed the \emph{structure} of the generation process (small model drafts, large model verifies) without changing the model itself---here, we change the structure of reasoning (parallel exploration, consensus-based routing) without changing the model or its training.
We show in \S\ref{sec:theory} that multi-path agreement provides an information-theoretically optimal difficulty estimator: when each path has independent error probability $\varepsilon$, the consensus signal concentrates exponentially in $K$, yielding a Spearman correlation with true difficulty of $\rho = -0.51$ ($p < 10^{-88}$)---far stronger than any single-path feature we tested.

\paragraph{The \method{} framework.}
We introduce \textbf{\fullmethod{}} (\method{}), a training-free, model-agnostic framework for adaptive test-time compute that operates in three phases:
\begin{enumerate}[nosep,leftmargin=*]
    \item \textbf{Explore.} Given a question $x$, generate $K$ short reasoning paths $\{r_1, \ldots, r_K\}$, each with a small token budget $b_{\text{explore}}$. These paths are independent and can be parallelized.
    \item \textbf{Decide.} Extract the answer from each path and compute a consensus score $c(x) = \max_a \frac{1}{K} \sum_{k=1}^{K} \mathbf{1}[a_k = a]$. Route the problem based on $c(x)$: high consensus $\rightarrow$ accept the majority answer immediately (easy); partial consensus $\rightarrow$ extend the most promising path with additional tokens (medium); low consensus $\rightarrow$ invoke full deliberation with all $K$ paths as context (hard).
    \item \textbf{Resolve.} For medium and hard problems, generate extended reasoning conditioned on the exploration paths, using a budget calibrated to the estimated difficulty tier.
\end{enumerate}
\method{} requires no training, no access to model internals, and no modification to the underlying LLM.
It operates entirely at the API level and is compatible with any reasoning model that supports temperature-based sampling.
The total compute budget is controlled by $K$, $b_{\text{explore}}$, and the per-tier continuation budgets, yielding a flexible accuracy--cost tradeoff.

\paragraph{Contributions.}
Our work makes the following contributions:
\begin{itemize}[nosep,leftmargin=*]
    \item \textbf{A new computation paradigm for adaptive test-time compute.}
    We introduce parallel exploration with consensus-based routing as an alternative to single-path budget allocation.
    Unlike prior adaptive methods that extract features from one reasoning trace, \method{} restructures the computation to produce a multi-path consensus signal that is fundamentally more informative about problem difficulty.

    \item \textbf{Information-theoretic justification.}
    We prove that multi-path consensus dominates any single-path difficulty feature in the black-box setting (\S\ref{sec:theory}).
    The consensus signal concentrates exponentially in $K$, explaining both why single-path controllers fail and why even modest $K$ (e.g., $K=5$) suffices.

    \item \textbf{Comprehensive empirical evaluation.}
    On GSM8K with Qwen3-8B, \method{} achieves \textbf{XX.X\%} accuracy at \textbf{XX\%} of the token cost of a fixed high-budget baseline, representing a \textbf{+X.X\,pp} improvement at matched compute and \textbf{XX\%} token savings at matched accuracy.
    We provide full ablations over $K$, $b_{\text{explore}}$, consensus thresholds, and continuation strategies.

    \item \textbf{First systematic negative results for black-box controllers.}
    We show that \emph{all} major categories of single-path adaptive controllers (honest features, uncertainty routing, MCTS refinement) fail to outperform fixed budgets, with degradations of $5$--$7$\,pp.
    These negative results are themselves a contribution, delimiting what is achievable without structural changes to the computation.

    \item \textbf{Training-free and model-agnostic.}
    \method{} requires no fine-tuning, no reward models, and no access to logits or hidden states.
    It works with any LLM API that supports standard generation and is immediately deployable.
\end{itemize}


Large language models equipped with chain-of-thought reasoning have achieved remarkable progress on complex tasks by ``thinking longer''---generating extended reasoning traces before committing to an answer.
This paradigm, exemplified by OpenAI o1~\cite{openai2024o1}, QwQ~\cite{qwq2024}, and DeepSeek-R1~\cite{deepseekr1}, frames test-time performance as a scaling problem: allocate more tokens, and accuracy improves.
Yet the gains from extended reasoning are far from uniform across problem instances.
Easy problems are \emph{overthought}: the model wastes hundreds of tokens re-deriving trivial facts or, worse, ``talks itself out of'' a correct intermediate answer through unnecessary self-revision.
Hard problems are \emph{underthought}: a fixed budget truncates reasoning before the model reaches the critical insight.
The fundamental mismatch between uniform compute allocation and heterogeneous problem difficulty leaves substantial accuracy and efficiency gains on the table.

\paragraph{The overthinking problem is pervasive.}
Through systematic experiments across GSM8K with Qwen3-8B, we quantify the extent of this mismatch.
When the reasoning budget is increased from a short to a long allocation, \textbf{XX.X\%} of samples exhibit \emph{accuracy degradation}---they were answered correctly at a lower budget but incorrectly at a higher one.
This is not an isolated edge case; it is a structured phenomenon affecting roughly one-third of all instances.
The ideal solution is clear in principle: allocate compute adaptively, spending little on easy problems and much on hard ones.
An oracle that selects the optimal budget per instance achieves \textbf{XX.X\%} accuracy---far above any fixed budget---demonstrating the enormous headroom for adaptive methods.

\paragraph{Existing adaptive approaches fail in the black-box setting.}
If the gains from adaptive allocation are so large, why has no practical method captured them?
We conducted what is, to our knowledge, the first comprehensive evaluation of black-box adaptive budget controllers for reasoning LLMs.
We tested all major categories of single-path difficulty signals:
(i)~\emph{honest features} extracted from a partial reasoning trace (token utilization, answer presence, lexical coherence), yielding controllers at $-6.5$\,pp below the fixed baseline;
(ii)~\emph{uncertainty-based routing} using answer entropy or self-reported confidence, at $-6.4$\,pp;
and (iii)~\emph{MCTS-style exploration--refinement}, at $-5.6$\,pp.
\textbf{All single-path controllers failed}, often performing \emph{worse} than the na\"ive fixed budget they were designed to improve.
The root cause is informational: a single reasoning trace is fundamentally impoverished as a difficulty signal.
The model's own confidence is poorly calibrated, token-level statistics are noisy, and partial traces reveal little about whether the \emph{final} answer will be correct.
In short, one witness is not enough to judge whether a problem is easy or hard.

\paragraph{Key insight: change the computation structure.}
Rather than trying to extract a weak difficulty signal from a single reasoning path, we propose to \emph{restructure the computation itself} to produce a fundamentally richer signal.
The core idea is simple: instead of running one long chain of thought, first run $K$ short, independent reasoning paths in parallel.
The \emph{cross-path consensus}---the degree to which these $K$ independent ``witnesses'' agree on an answer---provides a difficulty signal that is categorically stronger than anything extractable from a single trace.
Intuitively, if $K$ independent short explorations all converge to the same answer, the problem is almost certainly easy; if they scatter across many answers, it is hard and warrants extended deliberation.
This is analogous in spirit to how speculative decoding~\cite{leviathan2023fast} changed the \emph{structure} of the generation process (small model drafts, large model verifies) without changing the model itself---here, we change the structure of reasoning (parallel exploration, consensus-based routing) without changing the model or its training.
We show in \S\ref{sec:theory} that multi-path agreement provides an information-theoretically optimal difficulty estimator: when each path has independent error probability $\varepsilon$, the consensus signal concentrates exponentially in $K$, yielding a Spearman correlation with true difficulty of $\rho = -0.51$ ($p < 10^{-88}$)---far stronger than any single-path feature we tested.

\paragraph{The \method{} framework.}
We introduce \textbf{\fullmethod{}} (\method{}), a training-free, model-agnostic framework for adaptive test-time compute that operates in three phases:
\begin{enumerate}[nosep,leftmargin=*]
    \item \textbf{Explore.} Given a question $x$, generate $K$ short reasoning paths $\{r_1, \ldots, r_K\}$, each with a small token budget $b_{\text{explore}}$. These paths are independent and can be parallelized.
    \item \textbf{Decide.} Extract the answer from each path and compute a consensus score $c(x) = \max_a \frac{1}{K} \sum_{k=1}^{K} \mathbf{1}[a_k = a]$. Route the problem based on $c(x)$: high consensus $\rightarrow$ accept the majority answer immediately (easy); partial consensus $\rightarrow$ extend the most promising path with additional tokens (medium); low consensus $\rightarrow$ invoke full deliberation with all $K$ paths as context (hard).
    \item \textbf{Resolve.} For medium and hard problems, generate extended reasoning conditioned on the exploration paths, using a budget calibrated to the estimated difficulty tier.
\end{enumerate}
\method{} requires no training, no access to model internals, and no modification to the underlying LLM.
It operates entirely at the API level and is compatible with any reasoning model that supports temperature-based sampling.
The total compute budget is controlled by $K$, $b_{\text{explore}}$, and the per-tier continuation budgets, yielding a flexible accuracy--cost tradeoff.

\paragraph{Contributions.}
Our work makes the following contributions:
\begin{itemize}[nosep,leftmargin=*]
    \item \textbf{A new computation paradigm for adaptive test-time compute.}
    We introduce parallel exploration with consensus-based routing as an alternative to single-path budget allocation.
    Unlike prior adaptive methods that extract features from one reasoning trace, \method{} restructures the computation to produce a multi-path consensus signal that is fundamentally more informative about problem difficulty.

    \item \textbf{Information-theoretic justification.}
    We prove that multi-path consensus dominates any single-path difficulty feature in the black-box setting (\S\ref{sec:theory}).
    The consensus signal concentrates exponentially in $K$, explaining both why single-path controllers fail and why even modest $K$ (e.g., $K=5$) suffices.

    \item \textbf{Comprehensive empirical evaluation.}
    On GSM8K with Qwen3-8B, \method{} achieves \textbf{XX.X\%} accuracy at \textbf{XX\%} of the token cost of a fixed high-budget baseline, representing a \textbf{+X.X\,pp} improvement at matched compute and \textbf{XX\%} token savings at matched accuracy.
    We provide full ablations over $K$, $b_{\text{explore}}$, consensus thresholds, and continuation strategies.

    \item \textbf{First systematic negative results for black-box controllers.}
    We show that \emph{all} major categories of single-path adaptive controllers (honest features, uncertainty routing, MCTS refinement) fail to outperform fixed budgets, with degradations of $5$--$7$\,pp.
    These negative results are themselves a contribution, delimiting what is achievable without structural changes to the computation.

    \item \textbf{Training-free and model-agnostic.}
    \method{} requires no fine-tuning, no reward models, and no access to logits or hidden states.
    It works with any LLM API that supports standard generation and is immediately deployable.
\end{itemize}


% ============================================================================
\section{Related Work}
\label{sec:related}
% ============================================================================

% related_reasonspec.tex — Related Work for the Reasoning Speculation paper
% Loaded by main_reasonspec.tex via % related_reasonspec.tex — Related Work for the Reasoning Speculation paper
% Loaded by main_reasonspec.tex via % related_reasonspec.tex — Related Work for the Reasoning Speculation paper
% Loaded by main_reasonspec.tex via \input{sections/related_reasonspec}

% ---------------------------------------------------------------------------
\subsection{Test-Time Compute Scaling}
\label{sec:related:ttc}
% ---------------------------------------------------------------------------

The observation that ``thinking longer'' improves accuracy has motivated a new class of reasoning-specialized LLMs.
OpenAI o1~\citep{openai2024o1} first demonstrated that training models for extended chain-of-thought generation yields substantial gains on mathematical and scientific benchmarks.
QwQ~\citep{qwq2024} and DeepSeek-R1~\citep{deepseekr1} further showed that reinforcement-learning-based training for long reasoning produces strong open-weight reasoners.
\citet{snell2024scaling} formalized this as a \emph{test-time compute scaling law}: accuracy improves log-linearly with token budget on mathematical tasks.
However, this scaling is far from free---\citet{chen2024overthinking} and our own analysis (\S\ref{sec:intro}) show that accuracy \emph{degrades} on a non-trivial fraction of instances when the budget is increased, a phenomenon termed ``overthinking.''
Budget forcing~\citep{muennighoff2025s1} represents an early attempt at controlling reasoning length by truncating the thinking trace at a fixed token limit and appending a wrap-up prompt.
While effective as a diagnostic tool, budget forcing applies a uniform budget to all instances and thus cannot exploit the heterogeneity of problem difficulty---precisely the gap \method{} is designed to fill.
The ``sweet spot'' analysis of~\citet{chen2024overthinking} confirms that every instance has an optimal budget, but no fixed allocation can simultaneously serve easy and hard problems.

% ---------------------------------------------------------------------------
\subsection{Adaptive Budget Allocation}
\label{sec:related:adaptive}
% ---------------------------------------------------------------------------

A growing literature seeks to allocate reasoning compute \emph{per instance}, which we organize into three families.

\paragraph{Training-based methods.}
Several approaches modify the model or its training objective to produce difficulty-calibrated traces.
AdaThink~\citep{adathink2025} shapes rewards via GRPO to encourage the model to terminate reasoning early when confident, achieving budget savings without retraining from scratch.
AdaToken~\citep{adatoken2025} frames budget allocation as an online bandit problem and provides regret bounds, but requires online interaction during inference.
SelfBudgeter~\citep{liu2025selfbudgeter} trains the model to predict its own required budget before generation, enforcing adherence via budget-guided GRPO.
TALE~\citep{han2025tale} includes the target token budget in the prompt and trains via SFT/DPO to compress reasoning.
While these methods can be effective, they require model fine-tuning and are thus coupled to specific model weights and training data---a significant deployment constraint that \method{} avoids entirely.

\paragraph{Confidence-based methods.}
An alternative family uses the model's self-reported or estimated confidence to trigger early exit.
Satisficing Reasoning~\citep{luo2025satisficing} applies a confidence threshold to halt generation when the model deems its current answer sufficient.
Certaindex~\citep{luo2025certaindex} and Dynasor~\citep{fu2025dynasor} monitor uncertainty during generation and route to shorter or longer reasoning accordingly.
These methods are training-free but rely on a \emph{single-path} confidence signal, which we show is fundamentally limited: the model's self-assessed confidence is poorly calibrated, and single-trace uncertainty estimates achieve Spearman $|\rho| < 0.2$ with true problem difficulty (\S\ref{sec:theory}).

\paragraph{Training-free injection methods.}
RTB~\citep{hou2025thinkless} and Think Less Think Better~\citep{hou2025thinkless} inject wrap-up tokens (e.g., ``\texttt{Wait, let me reconsider\ldots}'') into the reasoning trace to steer the model toward shorter or longer deliberation without any training.
In our comprehensive evaluation (\S\ref{sec:intro}), we tested all major categories of single-path controllers---including feature-based, uncertainty-based, and injection-based methods---and found that \textbf{none outperformed the fixed-budget baseline}, with degradations of 5--7\,pp.
The root cause is informational: \emph{any} difficulty signal extracted from a single reasoning trace is bounded in mutual information with problem difficulty by $I(D; X_1)$ (Proposition~\ref{prop:single-path}), which is strictly less than the consensus signal $I(D; C_K)$ for $K \ge 2$.
This negative result motivates \method{}'s fundamental departure: rather than extracting a weak signal from one path, restructure the computation to produce a multi-path signal.

% ---------------------------------------------------------------------------
\subsection{Self-Consistency and Multi-Path Reasoning}
\label{sec:related:sc}
% ---------------------------------------------------------------------------

Self-consistency~\citep{wang2023selfconsistency} improves accuracy by generating $K$ complete reasoning traces and returning the majority-vote answer.
Best-of-$N$ with reward models~\citep{cobbe2021gsm8k,lightman2023prm} samples $N$ solutions and selects the highest-scoring one according to a learned verifier.
MCTS-based reasoning methods~\citep{hao2023reasoning,xie2024selfevaluation,qi2024mutual} build search trees over reasoning steps, using rollout-based value estimates to guide expansion.
These approaches all use multiple paths to improve \emph{answer quality}---but they allocate compute \emph{uniformly} across instances, spending the same $K$ full-budget traces on every question regardless of difficulty.
\method{} draws on the same insight that multi-path agreement correlates with correctness, but uses consensus as a \emph{difficulty signal for adaptive routing}, not merely as an aggregation mechanism.
Concretely, \method{} runs $K$ \emph{short} probes (not full-budget traces), uses the agreement pattern to \emph{decide} whether additional compute is warranted, and invests that compute \emph{selectively} on hard instances.
This yields a fundamentally different cost profile: self-consistency pays $K \times B_{\mathrm{full}}$ tokens per question; \method{} pays $K \times B_{\mathrm{probe}}$ for easy questions and scales up only where consensus indicates difficulty.

% ---------------------------------------------------------------------------
\subsection{Speculative Decoding}
\label{sec:related:spec}
% ---------------------------------------------------------------------------

Speculative decoding~\citep{leviathan2023fast,chen2023accelerating} accelerates autoregressive generation by using a small, fast draft model to propose token sequences that a larger target model then verifies in parallel.
The key insight is that \emph{restructuring the computation}---from sequential single-model generation to a draft-then-verify pipeline---can reduce latency without changing the output distribution.
This paradigm has been extended to tree-structured speculation~\citep{miao2024specinfer}, self-speculation within a single model~\citep{zhang2024draft}, and cascade routing across model sizes~\citep{chen2023frugalgpt}.

\method{} draws a deliberate analogy to speculative decoding, applied at the level of \emph{reasoning} rather than token generation.
Where speculative decoding uses a cheap draft model to identify tokens the target model would likely accept, \method{} uses cheap parallel probes to identify questions the model can likely answer without extended deliberation.
Where speculative decoding falls back to the target model for tokens the draft model gets wrong, \method{} falls back to full deliberation for questions where the probes disagree.
In both cases, the approach exploits a structural asymmetry---most tokens are easy to predict, most questions are easy to solve---to reduce expected cost while preserving or improving quality, without modifying the underlying model.


% ---------------------------------------------------------------------------
\subsection{Test-Time Compute Scaling}
\label{sec:related:ttc}
% ---------------------------------------------------------------------------

The observation that ``thinking longer'' improves accuracy has motivated a new class of reasoning-specialized LLMs.
OpenAI o1~\citep{openai2024o1} first demonstrated that training models for extended chain-of-thought generation yields substantial gains on mathematical and scientific benchmarks.
QwQ~\citep{qwq2024} and DeepSeek-R1~\citep{deepseekr1} further showed that reinforcement-learning-based training for long reasoning produces strong open-weight reasoners.
\citet{snell2024scaling} formalized this as a \emph{test-time compute scaling law}: accuracy improves log-linearly with token budget on mathematical tasks.
However, this scaling is far from free---\citet{chen2024overthinking} and our own analysis (\S\ref{sec:intro}) show that accuracy \emph{degrades} on a non-trivial fraction of instances when the budget is increased, a phenomenon termed ``overthinking.''
Budget forcing~\citep{muennighoff2025s1} represents an early attempt at controlling reasoning length by truncating the thinking trace at a fixed token limit and appending a wrap-up prompt.
While effective as a diagnostic tool, budget forcing applies a uniform budget to all instances and thus cannot exploit the heterogeneity of problem difficulty---precisely the gap \method{} is designed to fill.
The ``sweet spot'' analysis of~\citet{chen2024overthinking} confirms that every instance has an optimal budget, but no fixed allocation can simultaneously serve easy and hard problems.

% ---------------------------------------------------------------------------
\subsection{Adaptive Budget Allocation}
\label{sec:related:adaptive}
% ---------------------------------------------------------------------------

A growing literature seeks to allocate reasoning compute \emph{per instance}, which we organize into three families.

\paragraph{Training-based methods.}
Several approaches modify the model or its training objective to produce difficulty-calibrated traces.
AdaThink~\citep{adathink2025} shapes rewards via GRPO to encourage the model to terminate reasoning early when confident, achieving budget savings without retraining from scratch.
AdaToken~\citep{adatoken2025} frames budget allocation as an online bandit problem and provides regret bounds, but requires online interaction during inference.
SelfBudgeter~\citep{liu2025selfbudgeter} trains the model to predict its own required budget before generation, enforcing adherence via budget-guided GRPO.
TALE~\citep{han2025tale} includes the target token budget in the prompt and trains via SFT/DPO to compress reasoning.
While these methods can be effective, they require model fine-tuning and are thus coupled to specific model weights and training data---a significant deployment constraint that \method{} avoids entirely.

\paragraph{Confidence-based methods.}
An alternative family uses the model's self-reported or estimated confidence to trigger early exit.
Satisficing Reasoning~\citep{luo2025satisficing} applies a confidence threshold to halt generation when the model deems its current answer sufficient.
Certaindex~\citep{luo2025certaindex} and Dynasor~\citep{fu2025dynasor} monitor uncertainty during generation and route to shorter or longer reasoning accordingly.
These methods are training-free but rely on a \emph{single-path} confidence signal, which we show is fundamentally limited: the model's self-assessed confidence is poorly calibrated, and single-trace uncertainty estimates achieve Spearman $|\rho| < 0.2$ with true problem difficulty (\S\ref{sec:theory}).

\paragraph{Training-free injection methods.}
RTB~\citep{hou2025thinkless} and Think Less Think Better~\citep{hou2025thinkless} inject wrap-up tokens (e.g., ``\texttt{Wait, let me reconsider\ldots}'') into the reasoning trace to steer the model toward shorter or longer deliberation without any training.
In our comprehensive evaluation (\S\ref{sec:intro}), we tested all major categories of single-path controllers---including feature-based, uncertainty-based, and injection-based methods---and found that \textbf{none outperformed the fixed-budget baseline}, with degradations of 5--7\,pp.
The root cause is informational: \emph{any} difficulty signal extracted from a single reasoning trace is bounded in mutual information with problem difficulty by $I(D; X_1)$ (Proposition~\ref{prop:single-path}), which is strictly less than the consensus signal $I(D; C_K)$ for $K \ge 2$.
This negative result motivates \method{}'s fundamental departure: rather than extracting a weak signal from one path, restructure the computation to produce a multi-path signal.

% ---------------------------------------------------------------------------
\subsection{Self-Consistency and Multi-Path Reasoning}
\label{sec:related:sc}
% ---------------------------------------------------------------------------

Self-consistency~\citep{wang2023selfconsistency} improves accuracy by generating $K$ complete reasoning traces and returning the majority-vote answer.
Best-of-$N$ with reward models~\citep{cobbe2021gsm8k,lightman2023prm} samples $N$ solutions and selects the highest-scoring one according to a learned verifier.
MCTS-based reasoning methods~\citep{hao2023reasoning,xie2024selfevaluation,qi2024mutual} build search trees over reasoning steps, using rollout-based value estimates to guide expansion.
These approaches all use multiple paths to improve \emph{answer quality}---but they allocate compute \emph{uniformly} across instances, spending the same $K$ full-budget traces on every question regardless of difficulty.
\method{} draws on the same insight that multi-path agreement correlates with correctness, but uses consensus as a \emph{difficulty signal for adaptive routing}, not merely as an aggregation mechanism.
Concretely, \method{} runs $K$ \emph{short} probes (not full-budget traces), uses the agreement pattern to \emph{decide} whether additional compute is warranted, and invests that compute \emph{selectively} on hard instances.
This yields a fundamentally different cost profile: self-consistency pays $K \times B_{\mathrm{full}}$ tokens per question; \method{} pays $K \times B_{\mathrm{probe}}$ for easy questions and scales up only where consensus indicates difficulty.

% ---------------------------------------------------------------------------
\subsection{Speculative Decoding}
\label{sec:related:spec}
% ---------------------------------------------------------------------------

Speculative decoding~\citep{leviathan2023fast,chen2023accelerating} accelerates autoregressive generation by using a small, fast draft model to propose token sequences that a larger target model then verifies in parallel.
The key insight is that \emph{restructuring the computation}---from sequential single-model generation to a draft-then-verify pipeline---can reduce latency without changing the output distribution.
This paradigm has been extended to tree-structured speculation~\citep{miao2024specinfer}, self-speculation within a single model~\citep{zhang2024draft}, and cascade routing across model sizes~\citep{chen2023frugalgpt}.

\method{} draws a deliberate analogy to speculative decoding, applied at the level of \emph{reasoning} rather than token generation.
Where speculative decoding uses a cheap draft model to identify tokens the target model would likely accept, \method{} uses cheap parallel probes to identify questions the model can likely answer without extended deliberation.
Where speculative decoding falls back to the target model for tokens the draft model gets wrong, \method{} falls back to full deliberation for questions where the probes disagree.
In both cases, the approach exploits a structural asymmetry---most tokens are easy to predict, most questions are easy to solve---to reduce expected cost while preserving or improving quality, without modifying the underlying model.


% ---------------------------------------------------------------------------
\subsection{Test-Time Compute Scaling}
\label{sec:related:ttc}
% ---------------------------------------------------------------------------

The observation that ``thinking longer'' improves accuracy has motivated a new class of reasoning-specialized LLMs.
OpenAI o1~\citep{openai2024o1} first demonstrated that training models for extended chain-of-thought generation yields substantial gains on mathematical and scientific benchmarks.
QwQ~\citep{qwq2024} and DeepSeek-R1~\citep{deepseekr1} further showed that reinforcement-learning-based training for long reasoning produces strong open-weight reasoners.
\citet{snell2024scaling} formalized this as a \emph{test-time compute scaling law}: accuracy improves log-linearly with token budget on mathematical tasks.
However, this scaling is far from free---\citet{chen2024overthinking} and our own analysis (\S\ref{sec:intro}) show that accuracy \emph{degrades} on a non-trivial fraction of instances when the budget is increased, a phenomenon termed ``overthinking.''
Budget forcing~\citep{muennighoff2025s1} represents an early attempt at controlling reasoning length by truncating the thinking trace at a fixed token limit and appending a wrap-up prompt.
While effective as a diagnostic tool, budget forcing applies a uniform budget to all instances and thus cannot exploit the heterogeneity of problem difficulty---precisely the gap \method{} is designed to fill.
The ``sweet spot'' analysis of~\citet{chen2024overthinking} confirms that every instance has an optimal budget, but no fixed allocation can simultaneously serve easy and hard problems.

% ---------------------------------------------------------------------------
\subsection{Adaptive Budget Allocation}
\label{sec:related:adaptive}
% ---------------------------------------------------------------------------

A growing literature seeks to allocate reasoning compute \emph{per instance}, which we organize into three families.

\paragraph{Training-based methods.}
Several approaches modify the model or its training objective to produce difficulty-calibrated traces.
AdaThink~\citep{adathink2025} shapes rewards via GRPO to encourage the model to terminate reasoning early when confident, achieving budget savings without retraining from scratch.
AdaToken~\citep{adatoken2025} frames budget allocation as an online bandit problem and provides regret bounds, but requires online interaction during inference.
SelfBudgeter~\citep{liu2025selfbudgeter} trains the model to predict its own required budget before generation, enforcing adherence via budget-guided GRPO.
TALE~\citep{han2025tale} includes the target token budget in the prompt and trains via SFT/DPO to compress reasoning.
While these methods can be effective, they require model fine-tuning and are thus coupled to specific model weights and training data---a significant deployment constraint that \method{} avoids entirely.

\paragraph{Confidence-based methods.}
An alternative family uses the model's self-reported or estimated confidence to trigger early exit.
Satisficing Reasoning~\citep{luo2025satisficing} applies a confidence threshold to halt generation when the model deems its current answer sufficient.
Certaindex~\citep{luo2025certaindex} and Dynasor~\citep{fu2025dynasor} monitor uncertainty during generation and route to shorter or longer reasoning accordingly.
These methods are training-free but rely on a \emph{single-path} confidence signal, which we show is fundamentally limited: the model's self-assessed confidence is poorly calibrated, and single-trace uncertainty estimates achieve Spearman $|\rho| < 0.2$ with true problem difficulty (\S\ref{sec:theory}).

\paragraph{Training-free injection methods.}
RTB~\citep{hou2025thinkless} and Think Less Think Better~\citep{hou2025thinkless} inject wrap-up tokens (e.g., ``\texttt{Wait, let me reconsider\ldots}'') into the reasoning trace to steer the model toward shorter or longer deliberation without any training.
In our comprehensive evaluation (\S\ref{sec:intro}), we tested all major categories of single-path controllers---including feature-based, uncertainty-based, and injection-based methods---and found that \textbf{none outperformed the fixed-budget baseline}, with degradations of 5--7\,pp.
The root cause is informational: \emph{any} difficulty signal extracted from a single reasoning trace is bounded in mutual information with problem difficulty by $I(D; X_1)$ (Proposition~\ref{prop:single-path}), which is strictly less than the consensus signal $I(D; C_K)$ for $K \ge 2$.
This negative result motivates \method{}'s fundamental departure: rather than extracting a weak signal from one path, restructure the computation to produce a multi-path signal.

% ---------------------------------------------------------------------------
\subsection{Self-Consistency and Multi-Path Reasoning}
\label{sec:related:sc}
% ---------------------------------------------------------------------------

Self-consistency~\citep{wang2023selfconsistency} improves accuracy by generating $K$ complete reasoning traces and returning the majority-vote answer.
Best-of-$N$ with reward models~\citep{cobbe2021gsm8k,lightman2023prm} samples $N$ solutions and selects the highest-scoring one according to a learned verifier.
MCTS-based reasoning methods~\citep{hao2023reasoning,xie2024selfevaluation,qi2024mutual} build search trees over reasoning steps, using rollout-based value estimates to guide expansion.
These approaches all use multiple paths to improve \emph{answer quality}---but they allocate compute \emph{uniformly} across instances, spending the same $K$ full-budget traces on every question regardless of difficulty.
\method{} draws on the same insight that multi-path agreement correlates with correctness, but uses consensus as a \emph{difficulty signal for adaptive routing}, not merely as an aggregation mechanism.
Concretely, \method{} runs $K$ \emph{short} probes (not full-budget traces), uses the agreement pattern to \emph{decide} whether additional compute is warranted, and invests that compute \emph{selectively} on hard instances.
This yields a fundamentally different cost profile: self-consistency pays $K \times B_{\mathrm{full}}$ tokens per question; \method{} pays $K \times B_{\mathrm{probe}}$ for easy questions and scales up only where consensus indicates difficulty.

% ---------------------------------------------------------------------------
\subsection{Speculative Decoding}
\label{sec:related:spec}
% ---------------------------------------------------------------------------

Speculative decoding~\citep{leviathan2023fast,chen2023accelerating} accelerates autoregressive generation by using a small, fast draft model to propose token sequences that a larger target model then verifies in parallel.
The key insight is that \emph{restructuring the computation}---from sequential single-model generation to a draft-then-verify pipeline---can reduce latency without changing the output distribution.
This paradigm has been extended to tree-structured speculation~\citep{miao2024specinfer}, self-speculation within a single model~\citep{zhang2024draft}, and cascade routing across model sizes~\citep{chen2023frugalgpt}.

\method{} draws a deliberate analogy to speculative decoding, applied at the level of \emph{reasoning} rather than token generation.
Where speculative decoding uses a cheap draft model to identify tokens the target model would likely accept, \method{} uses cheap parallel probes to identify questions the model can likely answer without extended deliberation.
Where speculative decoding falls back to the target model for tokens the draft model gets wrong, \method{} falls back to full deliberation for questions where the probes disagree.
In both cases, the approach exploits a structural asymmetry---most tokens are easy to predict, most questions are easy to solve---to reduce expected cost while preserving or improving quality, without modifying the underlying model.


% ============================================================================
\section{Method: \fullmethod{}}
\label{sec:method}
% ============================================================================

% ===========================================================================
%  Method section for Reasoning Speculation (NeurIPS 2026)
% ===========================================================================

We now present \method{} (\fullmethod{}), a training-free, three-phase
framework that restructures test-time reasoning compute by using
\emph{parallel short explorations} to estimate problem difficulty via
cross-path consensus, then adaptively routing each question to an
appropriately sized resolution strategy.

% ---------------------------------------------------------------------------
\subsection{Problem Setup and Notation}
\label{sec:method:setup}
% ---------------------------------------------------------------------------

Let $\mathcal{M}$ denote a reasoning-capable language model (e.g.\ one with
a chain-of-thought or ``thinking'' mode) and let $q$ be an input question
drawn from a benchmark $\mathcal{Q}$.
A \emph{reasoning trace} $y = \mathcal{M}(q, b)$ is the text generated by
$\mathcal{M}$ when prompted with $q$ under a token budget~$b$.
We write $a(y)$ for the answer extracted from $y$ by a deterministic parser
(regex-based; see Appendix for details) and $|y|$ for the number of tokens
consumed.

We consider three budget levels
$B_{\mathrm{probe}} < B_{\mathrm{med}} < B_{\mathrm{hard}}$
(e.g.\ 128, 256, 512 tokens) and a parallelism factor $K \geq 2$.
The goal is an \emph{adaptive policy} $\pi$ that, given $q$ and
$\mathcal{M}$, selects an execution plan maximizing the cost-aware utility
\begin{equation}
  \label{eq:rs-utility}
  U(\pi) \;=\; \mathbb{E}_{q}\!\Big[\,
    \mathbf{1}\big[\,a(\pi(q)) = a^{\star}(q)\,\big]
    \;-\; \lambda \cdot \frac{|\pi(q)|}{B_{\mathrm{hard}}}
  \Big],
\end{equation}
where $a^{\star}(q)$ is the gold answer and $\lambda \geq 0$ trades off
accuracy against token cost.
A \emph{fixed-budget} baseline simply sets $\pi(q) = \mathcal{M}(q, B)$
for a constant~$B$; \emph{self-consistency}~(SC@$K$) samples $K$ full-budget
traces and majority-votes.
Both are compute-agnostic---they allocate the same resources regardless of
problem difficulty.

% ---------------------------------------------------------------------------
\subsection{Phase~1: Parallel Exploration}
\label{sec:method:explore}
% ---------------------------------------------------------------------------

Given question $q$, \method{} generates $K$ short \emph{exploration paths}
$\{y_1, \ldots, y_K\}$, each constrained to $B_{\mathrm{probe}}$ tokens:
\begin{equation}
  \label{eq:explore}
  y_k = \mathcal{M}(q,\, B_{\mathrm{probe}},\, \tau_k),
  \quad
  \tau_k =
  \begin{cases}
    0   & k = 1 \;\text{(greedy)}, \\
    \tau & k \geq 2 \;\text{(stochastic, } \tau > 0\text{)}.
  \end{cases}
\end{equation}
The first path uses greedy decoding ($\tau_1 = 0$) to obtain a high-quality
deterministic baseline; the remaining $K{-}1$ paths use temperature
$\tau > 0$ (default 0.7) to encourage diverse reasoning strategies.
This design mirrors the ``draft'' phase of speculative
decoding~\citep{leviathan2023fast,chen2023accelerating}: it produces
lightweight probes whose primary purpose is to \emph{estimate difficulty},
not to solve the problem outright.

\paragraph{Answer extraction and projection.}
From each path $y_k$, we extract a candidate answer $a_k = a(y_k)$ via
deterministic parsing.
Because $B_{\mathrm{probe}}$ may be insufficient for the model to reach a
``Final answer: \ldots'' marker, some paths yield $a_k = \varnothing$.
For these, we apply a lightweight \emph{projection} step: a short prompt
concatenates $q$ and the partial trace $y_k$ and asks the model to output a
single answer line, consuming at most 16 additional tokens.
Projection rescues partial reasoning that is substantively correct but
merely truncated before an answer marker, and adds negligible cost.

% ---------------------------------------------------------------------------
\subsection{Phase~2: Consensus-Based Difficulty Estimation}
\label{sec:method:decide}
% ---------------------------------------------------------------------------

The key insight underlying \method{} is that \emph{multi-path agreement is a
far more reliable difficulty signal than any feature extractable from a
single reasoning trace}.
Prior work~\citep{wang2023selfconsistency} observed that self-consistency
correlates with correctness; we operationalize this as a three-component
\emph{consensus signal} that drives an explicit routing decision.

\paragraph{Consensus signal.}
Let $\mathcal{A} = \{a_k : a_k \neq \varnothing\}$ be the set of valid
extracted answers and let $a^{\mathrm{maj}}$ be the answer with the highest
frequency in $\mathcal{A}$ (ties broken arbitrarily).
We define:
\begin{align}
  \label{eq:agreement}
  \alpha &= \frac{|\{a_k \in \mathcal{A} : a_k = a^{\mathrm{maj}}\}|}
                  {|\mathcal{A}|},
  &\quad&\text{(agreement ratio)}\\[4pt]
  \label{eq:validity}
  \nu &= \frac{|\mathcal{A}|}{K},
  &\quad&\text{(prediction validity)}\\[4pt]
  \label{eq:marker}
  \mu &= \frac{1}{K}\sum_{k=1}^{K}
         \mathbf{1}[\text{path } y_k \text{ contains a ``Final answer'' marker}],
  &\quad&\text{(final-marker rate)}
\end{align}
and an overall confidence score
\begin{equation}
  \label{eq:confidence}
  \gamma = 0.5\,\alpha + 0.3\,\nu + 0.2\,\mu \;\;\in [0, 1].
\end{equation}
Agreement $\alpha$ receives the largest weight because it directly measures
whether independent reasoning chains converge to the same answer.
Validity $\nu$ captures whether the probe budget was sufficient for the model
to produce any answer at all, and marker rate $\mu$ indicates whether traces
reached a natural conclusion.

\paragraph{Routing rule.}
The consensus signal determines the difficulty route:
\begin{equation}
  \label{eq:routing}
  \mathrm{route}(q) =
  \begin{cases}
    \textsc{Easy}
      & \text{if } \alpha \geq \theta_{\mathrm{E}} \;\text{and}\; \gamma \geq 0.6, \\[3pt]
    \textsc{Medium}
      & \text{if } \alpha \geq \theta_{\mathrm{M}} \;\text{or}\; \gamma \geq 0.4, \\[3pt]
    \textsc{Hard}
      & \text{otherwise},
  \end{cases}
\end{equation}
where $\theta_{\mathrm{E}}$ and $\theta_{\mathrm{M}}$ are agreement
thresholds (default 0.75 and 0.5, respectively).
The routing is deliberately conservative: a question is labeled \textsc{Easy}
only when a super-majority of paths agree \emph{and} overall confidence is
high; the \textsc{Hard} route is reserved for genuine disagreement.

% ---------------------------------------------------------------------------
\subsection{Phase~3: Adaptive Resolution}
\label{sec:method:resolve}
% ---------------------------------------------------------------------------

Each route prescribes a resolution strategy with a distinct token budget:

\paragraph{\textsc{Easy} route\,---\,majority vote.}
When paths exhibit strong consensus, the majority answer $a^{\mathrm{maj}}$
is returned immediately.
No additional generation is performed.

\paragraph{\textsc{Medium} route\,---\,extend the best path.}
Among exploration paths whose answer matches $a^{\mathrm{maj}}$, we select
the best path $y^{\star}$ (defaulting to the greedy path $y_1$) and
\emph{extend} its reasoning by continuing generation from $y^{\star}$ up to
$B_{\mathrm{med}}$ total tokens.
If the extension reveals a revised answer, we adopt it; otherwise, the
original answer from $y^{\star}$ is retained.
Concretely, the remaining budget is $B_{\mathrm{med}} - |y^{\star}|$
tokens, and extension uses greedy decoding ($\tau = 0$) to maximize
reliability.

\paragraph{\textsc{Hard} route\,---\,full deliberation.}
For problems with low consensus, we construct a \emph{deliberation prompt}
that includes the original question $q$ and summaries of all $K$ exploration
paths (truncated to avoid exceeding the context window).
The model then generates a fresh solution of up to $B_{\mathrm{hard}}$
tokens, informed by the diverse partial solutions as context:
\begin{equation}
  y_{\mathrm{delib}} = \mathcal{M}\!\big(\,
    [q;\; \mathrm{summarize}(y_1, \ldots, y_K)],\;
    B_{\mathrm{hard}},\; \tau = 0
  \big).
\end{equation}
This deliberation is analogous to an ``ensemble distillation'' within a
single model: the diverse reasoning strategies surfaced during exploration
serve as scaffolding for a more thorough solution.

% ---------------------------------------------------------------------------
\subsection{Full Algorithm}
\label{sec:method:algorithm}
% ---------------------------------------------------------------------------

Algorithm~\ref{alg:reasonspec} summarizes the complete \method{} pipeline.

\begin{algorithm}[t]
\caption{\method{}: Reasoning Speculation}
\label{alg:reasonspec}
\begin{algorithmic}[1]
\REQUIRE Question $q$, model $\mathcal{M}$, parallelism $K$,
  budgets $(B_{\mathrm{probe}}, B_{\mathrm{med}}, B_{\mathrm{hard}})$,
  thresholds $(\theta_{\mathrm{E}}, \theta_{\mathrm{M}})$,
  temperature $\tau$
\ENSURE Answer $\hat{a}$, tokens consumed $T$

\medskip
\STATE \textbf{// Phase 1: Parallel Exploration}
\FOR{$k = 1$ \textbf{to} $K$}
  \STATE $\tau_k \leftarrow 0$ if $k{=}1$ else $\tau$
    \hfill\COMMENT{Greedy first, diverse rest}
  \STATE $y_k \leftarrow \mathcal{M}(q,\; B_{\mathrm{probe}},\; \tau_k)$
  \STATE $a_k \leftarrow \textsc{Parse}(y_k)$
  \IF{$a_k = \varnothing$}
    \STATE $a_k \leftarrow \textsc{Project}(q, y_k)$
      \hfill\COMMENT{Lightweight projection, $\leq 16$ tokens}
  \ENDIF
\ENDFOR

\medskip
\STATE \textbf{// Phase 2: Consensus \& Routing}
\STATE Compute $\alpha, \nu, \mu, \gamma$ via
  Eqs.~\eqref{eq:agreement}--\eqref{eq:confidence}
\STATE $\mathrm{route} \leftarrow$ apply Eq.~\eqref{eq:routing}
\STATE $T \leftarrow \sum_{k=1}^{K} |y_k|$
  \hfill\COMMENT{Exploration cost}

\medskip
\STATE \textbf{// Phase 3: Adaptive Resolution}
\IF{$\mathrm{route} = \textsc{Easy}$}
  \STATE $\hat{a} \leftarrow a^{\mathrm{maj}}$
    \hfill\COMMENT{Return majority vote; no extra tokens}
\ELSIF{$\mathrm{route} = \textsc{Medium}$}
  \STATE $y^{\star} \leftarrow$ best path matching $a^{\mathrm{maj}}$
  \STATE $y_{\mathrm{ext}} \leftarrow \mathcal{M}(
    [q;\, y^{\star}],\; B_{\mathrm{med}} - |y^{\star}|,\; 0)$
  \STATE $\hat{a} \leftarrow \textsc{Parse}(y^{\star} \!\oplus\! y_{\mathrm{ext}})$;
    \;$T \leftarrow T + |y_{\mathrm{ext}}|$
\ELSE \hfill\COMMENT{$\textsc{Hard}$}
  \STATE $p \leftarrow [q;\; \textsc{Summarize}(y_1, \ldots, y_K)]$
  \STATE $y_{\mathrm{delib}} \leftarrow \mathcal{M}(p,\;
    B_{\mathrm{hard}},\; 0)$
  \STATE $\hat{a} \leftarrow \textsc{Parse}(y_{\mathrm{delib}})$;
    \;$T \leftarrow T + |y_{\mathrm{delib}}|$
\ENDIF
\RETURN $\hat{a},\; T$
\end{algorithmic}
\end{algorithm}

% ---------------------------------------------------------------------------
\subsection{Complexity Analysis}
\label{sec:method:complexity}
% ---------------------------------------------------------------------------

\paragraph{Token cost by route.}
Let $C_{\mathrm{explore}} = K \cdot B_{\mathrm{probe}}$ denote the
worst-case exploration cost (actual cost may be lower when paths terminate
early).
The total token budget consumed per question is:
\begin{equation}
  \label{eq:cost}
  T(q) =
  \begin{cases}
    C_{\mathrm{explore}}
      & \textsc{Easy}, \\[2pt]
    C_{\mathrm{explore}} + (B_{\mathrm{med}} - |y^{\star}|)
      & \textsc{Medium}, \\[2pt]
    C_{\mathrm{explore}} + B_{\mathrm{hard}}
      & \textsc{Hard}.
  \end{cases}
\end{equation}
With default parameters ($K{=}4$, $B_{\mathrm{probe}}{=}128$,
$B_{\mathrm{med}}{=}256$, $B_{\mathrm{hard}}{=}512$), the three routes
cost at most 512, 640, and 1024 tokens respectively.
If a fraction $p_E$ of questions are routed as \textsc{Easy}, $p_M$ as
\textsc{Medium}, and $p_H$ as \textsc{Hard}, the expected cost satisfies
\begin{equation}
  \label{eq:expected-cost}
  \mathbb{E}[T] = C_{\mathrm{explore}}
    + p_M \cdot \bar{E}_{\mathrm{med}}
    + p_H \cdot B_{\mathrm{hard}},
\end{equation}
where $\bar{E}_{\mathrm{med}}$ is the average extension length for
\textsc{Medium} questions.
When the majority of questions are easy (as is typical for well-studied
benchmarks), the expected cost approaches $K \cdot B_{\mathrm{probe}}$,
which can be substantially below a single full-budget pass at
$B_{\mathrm{hard}}$.

\paragraph{Latency.}
The $K$ exploration paths are \emph{embarrassingly parallel}: with
sufficient hardware, Phase~1 completes in the wall-clock time of a single
$B_{\mathrm{probe}}$-token generation, not $K$ times that.
This is a practical advantage over sequential methods such as progressive
halting or iterative refinement.
The resolution phase is a single sequential generation, so the total
wall-clock latency (with $K$-way parallelism) is bounded by
$B_{\mathrm{probe}} + \max(0, B_{\mathrm{med}}, B_{\mathrm{hard}})$
tokens, comparable to a single full-budget pass for \textsc{Medium}
and \textsc{Hard} questions, and strictly less for \textsc{Easy} ones.

% ---------------------------------------------------------------------------
\subsection{Comparison with Related Approaches}
\label{sec:method:comparison}
% ---------------------------------------------------------------------------

\method{} shares components with several existing paradigms but differs in
how it combines them.

\paragraph{Self-consistency (SC).}
SC@$K$~\citep{wang2023selfconsistency} generates $K$ full-budget
traces and majority-votes.
\method{} also uses majority voting, but \emph{only as a subroutine for
easy questions}; critically, it applies voting to \emph{short} probes
($B_{\mathrm{probe}}$ tokens) rather than full-length traces, and uses the
agreement signal to \emph{decide} whether full-length generation is
warranted.
SC pays $K \times B_{\mathrm{full}}$ tokens unconditionally; \method{}
pays $K \times B_{\mathrm{probe}}$ for easy questions and invests
additional compute only where consensus is low.
When the fraction of easy questions is non-trivial, this yields large
savings: for instance, with $B_{\mathrm{probe}} / B_{\mathrm{hard}} = 1/4$
and 50\% easy questions, \method{} uses roughly $3/8$ the tokens of SC@$K$
at comparable accuracy.

\paragraph{Best-of-$N$ (BoN) and reward-model reranking.}
BoN generates $N$ complete solutions and selects the one with the highest
reward model score~\citep{cobbe2021training,lightman2024lets}.
This requires (i)~$N$ full-budget generations and (ii)~a trained reward or
verifier model.
\method{} requires neither: it generates only short probes, needs no
external verifier, and uses consensus as a zero-shot quality signal.
When consensus is insufficient (\textsc{Hard} route), \method{} invests in
\emph{a single deeper generation} informed by all probes, rather than hoping
that one of $N$ independent samples happens to be correct.

\paragraph{Monte Carlo Tree Search (MCTS) for reasoning.}
MCTS-based methods~\citep{hao2023reasoning,xie2024selfevaluation} build a
search tree over reasoning steps, using a value function to guide
expansion.
While powerful, they require step-level value estimation (often from a
separate model), multiple sequential rollouts, and careful tuning of
exploration constants.
\method{} operates at the \emph{trace level} rather than the step level,
needs no value model, and makes a single routing decision after a fixed
number of parallel probes---resulting in a simpler, more predictable
compute profile.

\paragraph{Analogy to speculative decoding.}
The name \fullmethod{} is chosen deliberately: just as speculative
decoding~\citep{leviathan2023fast,chen2023accelerating} changes the
\emph{compute structure} of token generation (small-model draft $\to$
large-model verification) without altering the output distribution,
\method{} changes the \emph{compute structure of reasoning} (parallel
short probes $\to$ consensus $\to$ adaptive-depth resolution) without
modifying the model itself.
Both methods exploit a structural asymmetry---most tokens/problems are easy
and can be handled cheaply---to reduce the expected cost while preserving
or improving quality.


% ============================================================================
\section{Theoretical Analysis}
\label{sec:theory}
% ============================================================================

% theory_reasonspec.tex — Theoretical Analysis for ReasonSpec
% Included from main_reasonspec.tex via % theory_reasonspec.tex — Theoretical Analysis for ReasonSpec
% Included from main_reasonspec.tex via % theory_reasonspec.tex — Theoretical Analysis for ReasonSpec
% Included from main_reasonspec.tex via \input{sections/theory_reasonspec}

We now provide a theoretical foundation for the central design choice in
\method{}: using multi-path consensus rather than single-path features as
the difficulty signal for budget allocation.
Let $q$ denote a question, $D(q)\in\{\textrm{easy},\textrm{medium},\textrm{hard}\}$
its latent difficulty, $X_i$ the answer produced by the $i$-th independent
reasoning path, and
$C_K \triangleq \frac{1}{\binom{K}{2}}\sum_{i<j}\mathbf{1}[X_i = X_j]$
the pairwise agreement ratio among $K$ paths.

%% ===== Proposition 1 =====
\begin{proposition}[Consensus is more informative than a single path]
\label{prop:info}
For any $K \ge 2$,
\begin{equation}
  I\!\left(D;\, C_K\right) \;\ge\; I\!\left(D;\, X_1\right),
  \label{eq:mi-bound}
\end{equation}
where $I(\cdot\,;\cdot)$ denotes mutual information.
The inequality is strict whenever there exist two difficulty levels
$d, d'$ such that the answer distributions $P(X_1\!\mid\! D\!=\!d)$
and $P(X_1\!\mid\! D\!=\!d')$ differ.
\end{proposition}

\begin{proof}[Proof sketch]
By assumption, paths are conditionally i.i.d.\ given $q$ and hence given
$D(q)$.
Consider the joint sample $\mathbf{X}_K = (X_1, \dots, X_K)$.
By the data-processing inequality applied to the Markov chain
$D \to \mathbf{X}_K \to C_K$, we have
$I(D; C_K) \le I(D; \mathbf{X}_K)$.
In the reverse direction, $X_1$ is a deterministic function of
$\mathbf{X}_K$, so
$I(D; X_1) \le I(D; \mathbf{X}_K)$.
The key step is to observe that $C_K$ is a symmetric function of all $K$
samples and therefore captures \emph{inter-path} variation that a single
sample cannot.
Formally, $C_K$ is a function of the empirical distribution
$\hat{P}_K = \frac{1}{K}\sum_i \delta_{X_i}$, which is a sufficient
statistic for the family
$\{P(\mathbf{X}_K \!\mid\! D\!=\!d)\}_{d}$ by the factorization theorem
(since $P(\mathbf{X}_K \!\mid\! D) = \prod_i P(X_i \!\mid\! D)$
depends on $\mathbf{X}_K$ only through $\hat{P}_K$).
Sufficiency implies $I(D; \hat{P}_K) = I(D; \mathbf{X}_K)$,
and because $C_K$ is an injective function of $\hat{P}_K$ when the answer
space has at most $K$ distinct values, we obtain
$I(D; C_K) = I(D; \mathbf{X}_K) \ge I(D; X_1)$.
Strictness follows because $K \ge 2$ independent samples from distinct
distributions yield strictly more Fisher information and hence strictly
more mutual information than a single sample, provided
$P(X_1 \!\mid\! D\!=\!d) \ne P(X_1 \!\mid\! D\!=\!d')$ for some
$d \ne d'$.
A complete proof is given in Appendix~\ref{app:proofs}.
\end{proof}

%% ===== Proposition 2 =====
\begin{proposition}[Consensus as a monotone difficulty estimator]
\label{prop:monotone}
Let $p(q) = P(\text{correct}\!\mid\! q)$ denote the per-question
accuracy of a single path.
If paths are conditionally independent given $q$ and the answer space
is binary (correct/incorrect), then:
\begin{enumerate}[nosep,leftmargin=*]
  \item $\mathbb{E}[C_K \mid q]$ is strictly increasing in $p(q)$ for
        $p(q) > \tfrac{1}{2}$.
  \item As $K \to \infty$,
    \begin{equation}
      C_K \;\xrightarrow{\;a.s.\;}\; p(q)^{2} + \bigl(1 - p(q)\bigr)^{2},
      \label{eq:consensus-limit}
    \end{equation}
    which is a strictly decreasing function of difficulty
    (lower $p(q) \Rightarrow$ lower $C_\infty$) on $p(q)\in(\tfrac12,1]$.
  \item For finite $K$, $\mathrm{Var}(C_K \mid q) = O(1/K)$, so the
        difficulty estimate concentrates rapidly.
\end{enumerate}
Consequently, if $P(\textrm{correct}\!\mid\!\textrm{easy}) >
P(\textrm{correct}\!\mid\!\textrm{hard})$, then
$\mathbb{E}[C_K \mid \textrm{easy}] > \mathbb{E}[C_K \mid \textrm{hard}]$.
\end{proposition}

\begin{proof}[Proof sketch]
Under conditional independence and a binary answer space,
$C_K = \binom{K}{2}^{-1}\sum_{i<j}\mathbf{1}[X_i = X_j]$
has expectation
$\mathbb{E}[C_K \mid q] = p(q)^2 + (1-p(q))^2$.
Differentiating with respect to $p$:
$\frac{d}{dp}[p^2 + (1-p)^2] = 2(2p - 1) > 0$ for $p > \tfrac{1}{2}$.
Statement~(2) follows from the strong law of large numbers applied to
$\hat{P}_K$, and~(3) from the bounded-differences inequality applied
to $C_K$ (changing one path shifts $C_K$ by at most $2/(K-1)$).
\end{proof}

\begin{remark}
When the answer space is multiclass (e.g., numerical answers with many
possible values), the correct-answer probability $p(q)$ still dominates
$C_K$ because incorrect answers are typically dispersed across many
values.
Our experiments confirm this: on GSM8K, the average number of distinct
incorrect answers per question is $3.2$, ensuring that consensus is
driven primarily by agreement on the correct answer.
\end{remark}

%% ===== Proposition 3 =====
\begin{proposition}[Fundamental limitation of single-path signals]
\label{prop:single-path}
In the black-box setting (no access to logits, hidden states, or
internal model representations), the only observable from a single
reasoning path is the generated text $X_1$.
For any difficulty estimator $f\colon \mathcal{X} \to \mathbb{R}$
that is a deterministic function of $X_1$:
\begin{equation}
  I\!\left(D;\, f(X_1)\right) \;\le\; I\!\left(D;\, X_1\right)
  \;\le\; I\!\left(D;\, C_K\right) \quad \text{for } K \ge 2.
  \label{eq:single-path-bound}
\end{equation}
\end{proposition}

\begin{proof}[Proof sketch]
The first inequality is the data-processing inequality applied to
$D \to X_1 \to f(X_1)$.
The second follows from Proposition~\ref{prop:info}.
The bound is tight only in the degenerate case where a single sample
fully determines $D$---which requires the answer distributions under
different difficulty levels to have disjoint support, a condition that
fails in practice (easy and hard questions can produce the same answer
text).
\end{proof}

\paragraph{Empirical validation.}
We compute the Spearman rank correlation between candidate difficulty
signals and the oracle difficulty label (fraction of budgets on which
the model errs).
Single-path features---answer presence, token utilization, and reasoning
coherence---achieve $|\rho| \in [0.08, 0.19]$.
In contrast, the consensus signal $C_K$ with $K = 8$ achieves
$\rho = -0.51$ ($p < 10^{-88}$).
Multi-path consensus also yields a 72.8\% precision in identifying
easy-vs-hard questions, closing 73\% of the gap between single-path
heuristics and the oracle classifier.
These numbers corroborate Propositions~\ref{prop:info}--\ref{prop:single-path}:
consensus is not merely incrementally better; it provides a qualitatively
different signal.

%% ===== Proposition 4: Budget allocation =====
\begin{proposition}[Optimal budget allocation is a threshold policy on consensus]
\label{prop:threshold}
Consider a budget set $\mathcal{B} = \{b_\ell, b_m, b_h\}$ with
$b_\ell < b_m < b_h$, and define the per-question expected reward:
\begin{equation}
  R(q, b) \;=\; P\!\left(\textrm{correct} \mid q, b\right) - \lambda \cdot
  \frac{b}{b_h},
  \label{eq:reward}
\end{equation}
where $\lambda \ge 0$ penalizes token cost.
Suppose (i)~$P(\textrm{correct} \mid q, b)$ is non-decreasing in $b$
with diminishing returns (concavity), and
(ii)~the marginal accuracy gain of increasing the budget is
non-increasing in $p(q)$ (easy questions benefit less from extra budget).
Then the optimal policy $\pi^*(q) = \argmax_{b \in \mathcal{B}} R(q, b)$
takes the form of a \emph{threshold policy on $C_K$}:
\begin{equation}
  \pi^*(q) =
  \begin{cases}
    b_\ell & \text{if } C_K(q) \ge \tau_1, \\
    b_m    & \text{if } \tau_2 \le C_K(q) < \tau_1, \\
    b_h    & \text{if } C_K(q) < \tau_2,
  \end{cases}
  \label{eq:threshold-policy}
\end{equation}
for thresholds $0 \le \tau_2 \le \tau_1 \le 1$ that depend on $\lambda$
and the accuracy--budget curve.
\end{proposition}

\begin{proof}[Proof sketch]
Under assumptions~(i) and (ii), the marginal value of upgrading from
$b_\ell$ to $b_m$ (or $b_m$ to $b_h$) is a non-increasing function of
$p(q)$.
By Proposition~\ref{prop:monotone}, $C_K$ is a monotone proxy for
$p(q)$, so the indifference points where the marginal accuracy gain
equals the marginal cost $\lambda(b_{k+1}-b_k)/b_h$ correspond to
thresholds on $C_K$.
Between these thresholds, the optimal action is constant.
Concavity of the accuracy--budget curve ensures that at most three
regimes arise, matching the three routes in \method{}.
A formal derivation via Lagrangian relaxation is in
Appendix~\ref{app:proofs}.
\end{proof}

\paragraph{Connection to \method{} design.}
Proposition~\ref{prop:threshold} provides a normative justification for
the three-route architecture of \method{}:
high consensus ($C_K \ge \tau_1$) triggers the \textbf{Accept} route
(easy),
partial consensus ($\tau_2 \le C_K < \tau_1$) triggers the \textbf{Extend}
route (medium),
and low consensus ($C_K < \tau_2$) triggers the \textbf{Deliberate} route
(hard).
The thresholds $\tau_1, \tau_2$ are the only hyperparameters introduced by
\method{}, and they can be set by a lightweight sweep on a small
validation set or derived from the cost parameter $\lambda$.
Crucially, this result shows that no finer partition of the consensus
space is needed under the stated regularity conditions---the three-route
design is \emph{optimal} within the class of consensus-based policies.


We now provide a theoretical foundation for the central design choice in
\method{}: using multi-path consensus rather than single-path features as
the difficulty signal for budget allocation.
Let $q$ denote a question, $D(q)\in\{\textrm{easy},\textrm{medium},\textrm{hard}\}$
its latent difficulty, $X_i$ the answer produced by the $i$-th independent
reasoning path, and
$C_K \triangleq \frac{1}{\binom{K}{2}}\sum_{i<j}\mathbf{1}[X_i = X_j]$
the pairwise agreement ratio among $K$ paths.

%% ===== Proposition 1 =====
\begin{proposition}[Consensus is more informative than a single path]
\label{prop:info}
For any $K \ge 2$,
\begin{equation}
  I\!\left(D;\, C_K\right) \;\ge\; I\!\left(D;\, X_1\right),
  \label{eq:mi-bound}
\end{equation}
where $I(\cdot\,;\cdot)$ denotes mutual information.
The inequality is strict whenever there exist two difficulty levels
$d, d'$ such that the answer distributions $P(X_1\!\mid\! D\!=\!d)$
and $P(X_1\!\mid\! D\!=\!d')$ differ.
\end{proposition}

\begin{proof}[Proof sketch]
By assumption, paths are conditionally i.i.d.\ given $q$ and hence given
$D(q)$.
Consider the joint sample $\mathbf{X}_K = (X_1, \dots, X_K)$.
By the data-processing inequality applied to the Markov chain
$D \to \mathbf{X}_K \to C_K$, we have
$I(D; C_K) \le I(D; \mathbf{X}_K)$.
In the reverse direction, $X_1$ is a deterministic function of
$\mathbf{X}_K$, so
$I(D; X_1) \le I(D; \mathbf{X}_K)$.
The key step is to observe that $C_K$ is a symmetric function of all $K$
samples and therefore captures \emph{inter-path} variation that a single
sample cannot.
Formally, $C_K$ is a function of the empirical distribution
$\hat{P}_K = \frac{1}{K}\sum_i \delta_{X_i}$, which is a sufficient
statistic for the family
$\{P(\mathbf{X}_K \!\mid\! D\!=\!d)\}_{d}$ by the factorization theorem
(since $P(\mathbf{X}_K \!\mid\! D) = \prod_i P(X_i \!\mid\! D)$
depends on $\mathbf{X}_K$ only through $\hat{P}_K$).
Sufficiency implies $I(D; \hat{P}_K) = I(D; \mathbf{X}_K)$,
and because $C_K$ is an injective function of $\hat{P}_K$ when the answer
space has at most $K$ distinct values, we obtain
$I(D; C_K) = I(D; \mathbf{X}_K) \ge I(D; X_1)$.
Strictness follows because $K \ge 2$ independent samples from distinct
distributions yield strictly more Fisher information and hence strictly
more mutual information than a single sample, provided
$P(X_1 \!\mid\! D\!=\!d) \ne P(X_1 \!\mid\! D\!=\!d')$ for some
$d \ne d'$.
A complete proof is given in Appendix~\ref{app:proofs}.
\end{proof}

%% ===== Proposition 2 =====
\begin{proposition}[Consensus as a monotone difficulty estimator]
\label{prop:monotone}
Let $p(q) = P(\text{correct}\!\mid\! q)$ denote the per-question
accuracy of a single path.
If paths are conditionally independent given $q$ and the answer space
is binary (correct/incorrect), then:
\begin{enumerate}[nosep,leftmargin=*]
  \item $\mathbb{E}[C_K \mid q]$ is strictly increasing in $p(q)$ for
        $p(q) > \tfrac{1}{2}$.
  \item As $K \to \infty$,
    \begin{equation}
      C_K \;\xrightarrow{\;a.s.\;}\; p(q)^{2} + \bigl(1 - p(q)\bigr)^{2},
      \label{eq:consensus-limit}
    \end{equation}
    which is a strictly decreasing function of difficulty
    (lower $p(q) \Rightarrow$ lower $C_\infty$) on $p(q)\in(\tfrac12,1]$.
  \item For finite $K$, $\mathrm{Var}(C_K \mid q) = O(1/K)$, so the
        difficulty estimate concentrates rapidly.
\end{enumerate}
Consequently, if $P(\textrm{correct}\!\mid\!\textrm{easy}) >
P(\textrm{correct}\!\mid\!\textrm{hard})$, then
$\mathbb{E}[C_K \mid \textrm{easy}] > \mathbb{E}[C_K \mid \textrm{hard}]$.
\end{proposition}

\begin{proof}[Proof sketch]
Under conditional independence and a binary answer space,
$C_K = \binom{K}{2}^{-1}\sum_{i<j}\mathbf{1}[X_i = X_j]$
has expectation
$\mathbb{E}[C_K \mid q] = p(q)^2 + (1-p(q))^2$.
Differentiating with respect to $p$:
$\frac{d}{dp}[p^2 + (1-p)^2] = 2(2p - 1) > 0$ for $p > \tfrac{1}{2}$.
Statement~(2) follows from the strong law of large numbers applied to
$\hat{P}_K$, and~(3) from the bounded-differences inequality applied
to $C_K$ (changing one path shifts $C_K$ by at most $2/(K-1)$).
\end{proof}

\begin{remark}
When the answer space is multiclass (e.g., numerical answers with many
possible values), the correct-answer probability $p(q)$ still dominates
$C_K$ because incorrect answers are typically dispersed across many
values.
Our experiments confirm this: on GSM8K, the average number of distinct
incorrect answers per question is $3.2$, ensuring that consensus is
driven primarily by agreement on the correct answer.
\end{remark}

%% ===== Proposition 3 =====
\begin{proposition}[Fundamental limitation of single-path signals]
\label{prop:single-path}
In the black-box setting (no access to logits, hidden states, or
internal model representations), the only observable from a single
reasoning path is the generated text $X_1$.
For any difficulty estimator $f\colon \mathcal{X} \to \mathbb{R}$
that is a deterministic function of $X_1$:
\begin{equation}
  I\!\left(D;\, f(X_1)\right) \;\le\; I\!\left(D;\, X_1\right)
  \;\le\; I\!\left(D;\, C_K\right) \quad \text{for } K \ge 2.
  \label{eq:single-path-bound}
\end{equation}
\end{proposition}

\begin{proof}[Proof sketch]
The first inequality is the data-processing inequality applied to
$D \to X_1 \to f(X_1)$.
The second follows from Proposition~\ref{prop:info}.
The bound is tight only in the degenerate case where a single sample
fully determines $D$---which requires the answer distributions under
different difficulty levels to have disjoint support, a condition that
fails in practice (easy and hard questions can produce the same answer
text).
\end{proof}

\paragraph{Empirical validation.}
We compute the Spearman rank correlation between candidate difficulty
signals and the oracle difficulty label (fraction of budgets on which
the model errs).
Single-path features---answer presence, token utilization, and reasoning
coherence---achieve $|\rho| \in [0.08, 0.19]$.
In contrast, the consensus signal $C_K$ with $K = 8$ achieves
$\rho = -0.51$ ($p < 10^{-88}$).
Multi-path consensus also yields a 72.8\% precision in identifying
easy-vs-hard questions, closing 73\% of the gap between single-path
heuristics and the oracle classifier.
These numbers corroborate Propositions~\ref{prop:info}--\ref{prop:single-path}:
consensus is not merely incrementally better; it provides a qualitatively
different signal.

%% ===== Proposition 4: Budget allocation =====
\begin{proposition}[Optimal budget allocation is a threshold policy on consensus]
\label{prop:threshold}
Consider a budget set $\mathcal{B} = \{b_\ell, b_m, b_h\}$ with
$b_\ell < b_m < b_h$, and define the per-question expected reward:
\begin{equation}
  R(q, b) \;=\; P\!\left(\textrm{correct} \mid q, b\right) - \lambda \cdot
  \frac{b}{b_h},
  \label{eq:reward}
\end{equation}
where $\lambda \ge 0$ penalizes token cost.
Suppose (i)~$P(\textrm{correct} \mid q, b)$ is non-decreasing in $b$
with diminishing returns (concavity), and
(ii)~the marginal accuracy gain of increasing the budget is
non-increasing in $p(q)$ (easy questions benefit less from extra budget).
Then the optimal policy $\pi^*(q) = \argmax_{b \in \mathcal{B}} R(q, b)$
takes the form of a \emph{threshold policy on $C_K$}:
\begin{equation}
  \pi^*(q) =
  \begin{cases}
    b_\ell & \text{if } C_K(q) \ge \tau_1, \\
    b_m    & \text{if } \tau_2 \le C_K(q) < \tau_1, \\
    b_h    & \text{if } C_K(q) < \tau_2,
  \end{cases}
  \label{eq:threshold-policy}
\end{equation}
for thresholds $0 \le \tau_2 \le \tau_1 \le 1$ that depend on $\lambda$
and the accuracy--budget curve.
\end{proposition}

\begin{proof}[Proof sketch]
Under assumptions~(i) and (ii), the marginal value of upgrading from
$b_\ell$ to $b_m$ (or $b_m$ to $b_h$) is a non-increasing function of
$p(q)$.
By Proposition~\ref{prop:monotone}, $C_K$ is a monotone proxy for
$p(q)$, so the indifference points where the marginal accuracy gain
equals the marginal cost $\lambda(b_{k+1}-b_k)/b_h$ correspond to
thresholds on $C_K$.
Between these thresholds, the optimal action is constant.
Concavity of the accuracy--budget curve ensures that at most three
regimes arise, matching the three routes in \method{}.
A formal derivation via Lagrangian relaxation is in
Appendix~\ref{app:proofs}.
\end{proof}

\paragraph{Connection to \method{} design.}
Proposition~\ref{prop:threshold} provides a normative justification for
the three-route architecture of \method{}:
high consensus ($C_K \ge \tau_1$) triggers the \textbf{Accept} route
(easy),
partial consensus ($\tau_2 \le C_K < \tau_1$) triggers the \textbf{Extend}
route (medium),
and low consensus ($C_K < \tau_2$) triggers the \textbf{Deliberate} route
(hard).
The thresholds $\tau_1, \tau_2$ are the only hyperparameters introduced by
\method{}, and they can be set by a lightweight sweep on a small
validation set or derived from the cost parameter $\lambda$.
Crucially, this result shows that no finer partition of the consensus
space is needed under the stated regularity conditions---the three-route
design is \emph{optimal} within the class of consensus-based policies.


We now provide a theoretical foundation for the central design choice in
\method{}: using multi-path consensus rather than single-path features as
the difficulty signal for budget allocation.
Let $q$ denote a question, $D(q)\in\{\textrm{easy},\textrm{medium},\textrm{hard}\}$
its latent difficulty, $X_i$ the answer produced by the $i$-th independent
reasoning path, and
$C_K \triangleq \frac{1}{\binom{K}{2}}\sum_{i<j}\mathbf{1}[X_i = X_j]$
the pairwise agreement ratio among $K$ paths.

%% ===== Proposition 1 =====
\begin{proposition}[Consensus is more informative than a single path]
\label{prop:info}
For any $K \ge 2$,
\begin{equation}
  I\!\left(D;\, C_K\right) \;\ge\; I\!\left(D;\, X_1\right),
  \label{eq:mi-bound}
\end{equation}
where $I(\cdot\,;\cdot)$ denotes mutual information.
The inequality is strict whenever there exist two difficulty levels
$d, d'$ such that the answer distributions $P(X_1\!\mid\! D\!=\!d)$
and $P(X_1\!\mid\! D\!=\!d')$ differ.
\end{proposition}

\begin{proof}[Proof sketch]
By assumption, paths are conditionally i.i.d.\ given $q$ and hence given
$D(q)$.
Consider the joint sample $\mathbf{X}_K = (X_1, \dots, X_K)$.
By the data-processing inequality applied to the Markov chain
$D \to \mathbf{X}_K \to C_K$, we have
$I(D; C_K) \le I(D; \mathbf{X}_K)$.
In the reverse direction, $X_1$ is a deterministic function of
$\mathbf{X}_K$, so
$I(D; X_1) \le I(D; \mathbf{X}_K)$.
The key step is to observe that $C_K$ is a symmetric function of all $K$
samples and therefore captures \emph{inter-path} variation that a single
sample cannot.
Formally, $C_K$ is a function of the empirical distribution
$\hat{P}_K = \frac{1}{K}\sum_i \delta_{X_i}$, which is a sufficient
statistic for the family
$\{P(\mathbf{X}_K \!\mid\! D\!=\!d)\}_{d}$ by the factorization theorem
(since $P(\mathbf{X}_K \!\mid\! D) = \prod_i P(X_i \!\mid\! D)$
depends on $\mathbf{X}_K$ only through $\hat{P}_K$).
Sufficiency implies $I(D; \hat{P}_K) = I(D; \mathbf{X}_K)$,
and because $C_K$ is an injective function of $\hat{P}_K$ when the answer
space has at most $K$ distinct values, we obtain
$I(D; C_K) = I(D; \mathbf{X}_K) \ge I(D; X_1)$.
Strictness follows because $K \ge 2$ independent samples from distinct
distributions yield strictly more Fisher information and hence strictly
more mutual information than a single sample, provided
$P(X_1 \!\mid\! D\!=\!d) \ne P(X_1 \!\mid\! D\!=\!d')$ for some
$d \ne d'$.
A complete proof is given in Appendix~\ref{app:proofs}.
\end{proof}

%% ===== Proposition 2 =====
\begin{proposition}[Consensus as a monotone difficulty estimator]
\label{prop:monotone}
Let $p(q) = P(\text{correct}\!\mid\! q)$ denote the per-question
accuracy of a single path.
If paths are conditionally independent given $q$ and the answer space
is binary (correct/incorrect), then:
\begin{enumerate}[nosep,leftmargin=*]
  \item $\mathbb{E}[C_K \mid q]$ is strictly increasing in $p(q)$ for
        $p(q) > \tfrac{1}{2}$.
  \item As $K \to \infty$,
    \begin{equation}
      C_K \;\xrightarrow{\;a.s.\;}\; p(q)^{2} + \bigl(1 - p(q)\bigr)^{2},
      \label{eq:consensus-limit}
    \end{equation}
    which is a strictly decreasing function of difficulty
    (lower $p(q) \Rightarrow$ lower $C_\infty$) on $p(q)\in(\tfrac12,1]$.
  \item For finite $K$, $\mathrm{Var}(C_K \mid q) = O(1/K)$, so the
        difficulty estimate concentrates rapidly.
\end{enumerate}
Consequently, if $P(\textrm{correct}\!\mid\!\textrm{easy}) >
P(\textrm{correct}\!\mid\!\textrm{hard})$, then
$\mathbb{E}[C_K \mid \textrm{easy}] > \mathbb{E}[C_K \mid \textrm{hard}]$.
\end{proposition}

\begin{proof}[Proof sketch]
Under conditional independence and a binary answer space,
$C_K = \binom{K}{2}^{-1}\sum_{i<j}\mathbf{1}[X_i = X_j]$
has expectation
$\mathbb{E}[C_K \mid q] = p(q)^2 + (1-p(q))^2$.
Differentiating with respect to $p$:
$\frac{d}{dp}[p^2 + (1-p)^2] = 2(2p - 1) > 0$ for $p > \tfrac{1}{2}$.
Statement~(2) follows from the strong law of large numbers applied to
$\hat{P}_K$, and~(3) from the bounded-differences inequality applied
to $C_K$ (changing one path shifts $C_K$ by at most $2/(K-1)$).
\end{proof}

\begin{remark}
When the answer space is multiclass (e.g., numerical answers with many
possible values), the correct-answer probability $p(q)$ still dominates
$C_K$ because incorrect answers are typically dispersed across many
values.
Our experiments confirm this: on GSM8K, the average number of distinct
incorrect answers per question is $3.2$, ensuring that consensus is
driven primarily by agreement on the correct answer.
\end{remark}

%% ===== Proposition 3 =====
\begin{proposition}[Fundamental limitation of single-path signals]
\label{prop:single-path}
In the black-box setting (no access to logits, hidden states, or
internal model representations), the only observable from a single
reasoning path is the generated text $X_1$.
For any difficulty estimator $f\colon \mathcal{X} \to \mathbb{R}$
that is a deterministic function of $X_1$:
\begin{equation}
  I\!\left(D;\, f(X_1)\right) \;\le\; I\!\left(D;\, X_1\right)
  \;\le\; I\!\left(D;\, C_K\right) \quad \text{for } K \ge 2.
  \label{eq:single-path-bound}
\end{equation}
\end{proposition}

\begin{proof}[Proof sketch]
The first inequality is the data-processing inequality applied to
$D \to X_1 \to f(X_1)$.
The second follows from Proposition~\ref{prop:info}.
The bound is tight only in the degenerate case where a single sample
fully determines $D$---which requires the answer distributions under
different difficulty levels to have disjoint support, a condition that
fails in practice (easy and hard questions can produce the same answer
text).
\end{proof}

\paragraph{Empirical validation.}
We compute the Spearman rank correlation between candidate difficulty
signals and the oracle difficulty label (fraction of budgets on which
the model errs).
Single-path features---answer presence, token utilization, and reasoning
coherence---achieve $|\rho| \in [0.08, 0.19]$.
In contrast, the consensus signal $C_K$ with $K = 8$ achieves
$\rho = -0.51$ ($p < 10^{-88}$).
Multi-path consensus also yields a 72.8\% precision in identifying
easy-vs-hard questions, closing 73\% of the gap between single-path
heuristics and the oracle classifier.
These numbers corroborate Propositions~\ref{prop:info}--\ref{prop:single-path}:
consensus is not merely incrementally better; it provides a qualitatively
different signal.

%% ===== Proposition 4: Budget allocation =====
\begin{proposition}[Optimal budget allocation is a threshold policy on consensus]
\label{prop:threshold}
Consider a budget set $\mathcal{B} = \{b_\ell, b_m, b_h\}$ with
$b_\ell < b_m < b_h$, and define the per-question expected reward:
\begin{equation}
  R(q, b) \;=\; P\!\left(\textrm{correct} \mid q, b\right) - \lambda \cdot
  \frac{b}{b_h},
  \label{eq:reward}
\end{equation}
where $\lambda \ge 0$ penalizes token cost.
Suppose (i)~$P(\textrm{correct} \mid q, b)$ is non-decreasing in $b$
with diminishing returns (concavity), and
(ii)~the marginal accuracy gain of increasing the budget is
non-increasing in $p(q)$ (easy questions benefit less from extra budget).
Then the optimal policy $\pi^*(q) = \argmax_{b \in \mathcal{B}} R(q, b)$
takes the form of a \emph{threshold policy on $C_K$}:
\begin{equation}
  \pi^*(q) =
  \begin{cases}
    b_\ell & \text{if } C_K(q) \ge \tau_1, \\
    b_m    & \text{if } \tau_2 \le C_K(q) < \tau_1, \\
    b_h    & \text{if } C_K(q) < \tau_2,
  \end{cases}
  \label{eq:threshold-policy}
\end{equation}
for thresholds $0 \le \tau_2 \le \tau_1 \le 1$ that depend on $\lambda$
and the accuracy--budget curve.
\end{proposition}

\begin{proof}[Proof sketch]
Under assumptions~(i) and (ii), the marginal value of upgrading from
$b_\ell$ to $b_m$ (or $b_m$ to $b_h$) is a non-increasing function of
$p(q)$.
By Proposition~\ref{prop:monotone}, $C_K$ is a monotone proxy for
$p(q)$, so the indifference points where the marginal accuracy gain
equals the marginal cost $\lambda(b_{k+1}-b_k)/b_h$ correspond to
thresholds on $C_K$.
Between these thresholds, the optimal action is constant.
Concavity of the accuracy--budget curve ensures that at most three
regimes arise, matching the three routes in \method{}.
A formal derivation via Lagrangian relaxation is in
Appendix~\ref{app:proofs}.
\end{proof}

\paragraph{Connection to \method{} design.}
Proposition~\ref{prop:threshold} provides a normative justification for
the three-route architecture of \method{}:
high consensus ($C_K \ge \tau_1$) triggers the \textbf{Accept} route
(easy),
partial consensus ($\tau_2 \le C_K < \tau_1$) triggers the \textbf{Extend}
route (medium),
and low consensus ($C_K < \tau_2$) triggers the \textbf{Deliberate} route
(hard).
The thresholds $\tau_1, \tau_2$ are the only hyperparameters introduced by
\method{}, and they can be set by a lightweight sweep on a small
validation set or derived from the cost parameter $\lambda$.
Crucially, this result shows that no finer partition of the consensus
space is needed under the stated regularity conditions---the three-route
design is \emph{optimal} within the class of consensus-based policies.


% ============================================================================
\section{Experiments}
\label{sec:experiments}
% ============================================================================

% experiments_reasonspec.tex — Experiments section for Reasoning Speculation
% Included from main_reasonspec.tex via % experiments_reasonspec.tex — Experiments section for Reasoning Speculation
% Included from main_reasonspec.tex via % experiments_reasonspec.tex — Experiments section for Reasoning Speculation
% Included from main_reasonspec.tex via \input{sections/experiments_reasonspec}

% ---------------------------------------------------------------------------
\subsection{Experimental Setup}
\label{sec:exp:setup}
% ---------------------------------------------------------------------------

\paragraph{Model and inference.}
All experiments use \textbf{Qwen3-8B} with thinking mode enabled
(\texttt{enable\_thinking=True}), which instructs the model to produce an
extended chain-of-thought within \texttt{<think>...</think>} tags before
emitting a final answer.
We control the thinking token budget $b$ by setting
\texttt{max\_thinking\_tokens}; once this limit is reached, the model is
forced to exit the thinking block and produce an answer.
Greedy decoding ($\tau = 0$) is used for the first exploration path and
all resolution phases; the remaining $K{-}1$ exploration paths use
$\tau = 0.7$ to encourage diversity (\S\ref{sec:method:explore}).

\paragraph{Benchmark.}
We evaluate on \textbf{GSM8K}~\citep{cobbe2021training}, a dataset of
1{,}319 grade-school math word problems.
For the main experiments and ablations we sample subsets of $n = 200$
problems (with fixed random seeds for reproducibility); full-dataset
results are reported for the best configuration.
Accuracy is measured via exact match after deterministic number extraction
(the rightmost numerical value in the model's output).

\paragraph{Baselines.}
We compare \method{} against four baselines that span the natural
design space:
\begin{itemize}[nosep,leftmargin=*]
  \item \textbf{Fixed-$b$}: single-pass generation with a fixed thinking
    budget $b \in \{128, 256, 512\}$ tokens.
  \item \textbf{SC@$K{\times}128$}: self-consistency~\citep{wang2023selfconsistency}
    with $K$ independent samples at $b = 128$ tokens each, followed by
    majority voting.
    This baseline uses the same total exploration tokens as \method{}
    with $K$ paths and $B_{\mathrm{probe}} = 128$, providing a
    compute-matched comparison.
\end{itemize}

\paragraph{\method{} configurations.}
We explore the following hyperparameter axes:
parallelism $K \in \{2, 3, 4, 8\}$,
probe budget $B_{\mathrm{probe}} \in \{128, 170, 256, 512\}$,
medium resolution budget $B_{\mathrm{med}} \in \{256, 512\}$, and
hard resolution budget $B_{\mathrm{hard}} \in \{512, 1024\}$.
Routing thresholds default to $\theta_E = 0.75$ (easy) and
$\theta_M = 0.5$ (medium) unless varied in ablation studies.
The consensus confidence weights are fixed at
$\gamma = 0.5\alpha + 0.3\nu + 0.2\mu$ throughout.

\paragraph{Evaluation metrics.}
For each method we report:
(1)~\textbf{Accuracy} (exact-match \%),
(2)~\textbf{Avg Tokens} (mean thinking tokens consumed per sample,
including exploration overhead for \method{}),
(3)~\textbf{Route distribution} (fraction of samples classified as
Easy / Medium / Hard by the consensus router).

\paragraph{Hardware.}
All experiments are conducted on a single NVIDIA A100-80GB GPU.

% ---------------------------------------------------------------------------
\subsection{Main Results}
\label{sec:exp:main}
% ---------------------------------------------------------------------------

Table~\ref{tab:main-rs} presents the primary comparison between
\method{}, fixed-budget baselines, and self-consistency on GSM8K
($n = 200$).
All \method{} variants use default routing thresholds
($\theta_E = 0.75$, $\theta_M = 0.5$).

\begin{table}[t]
\centering
\caption{%
  \textbf{Main results on GSM8K} ($n = 200$, Qwen3-8B thinking mode).
  \textbf{Avg Tok} includes full end-to-end cost (exploration +
  resolution).
  Route columns report the fraction of samples classified as Easy (E),
  Medium (M), and Hard (H) by the consensus router.
  Best adaptive result in \textbf{bold}; best baseline \underline{underlined}.
}
\label{tab:main-rs}
\small
\setlength{\tabcolsep}{4pt}
\begin{tabular}{@{}lcccccccc@{}}
\toprule
\textbf{Method} & $K$ & $B_{\mathrm{probe}}$ &
  \textbf{Acc (\%)} & \textbf{Avg Tok} &
  \textbf{E\%} & \textbf{M\%} & \textbf{H\%} \\
\midrule
\multicolumn{8}{@{}l}{\textit{Fixed-budget baselines}} \\
Fixed-128       & -- & -- & XX.X & 128  & -- & -- & -- \\
Fixed-256       & -- & -- & XX.X & 256  & -- & -- & -- \\
\underline{Fixed-512} & -- & -- & \underline{XX.X} & 512  & -- & -- & -- \\
\midrule
\multicolumn{8}{@{}l}{\textit{Self-consistency baselines}} \\
SC@$2{\times}128$  & 2 & 128 & XX.X & 256  & -- & -- & -- \\
SC@$4{\times}128$  & 4 & 128 & XX.X & 512  & -- & -- & -- \\
SC@$8{\times}128$  & 8 & 128 & XX.X & 1024 & -- & -- & -- \\
\midrule
\multicolumn{8}{@{}l}{\textit{\method{} (ours)}} \\
\method{}-2      & 2 & 128 & XX.X & XX.X & XX.X & XX.X & XX.X \\
\method{}-4      & 4 & 128 & XX.X & XX.X & XX.X & XX.X & XX.X \\
\method{}-4-256  & 4 & 256 & XX.X & XX.X & XX.X & XX.X & XX.X \\
\textbf{\method{}-8} & 8 & 128 & \textbf{XX.X} & XX.X & XX.X & XX.X & XX.X \\
\method{}-8-256  & 8 & 256 & XX.X & XX.X & XX.X & XX.X & XX.X \\
\bottomrule
\end{tabular}
\end{table}

Several observations emerge from Table~\ref{tab:main-rs}.
First, \method{} consistently outperforms fixed-budget baselines at
comparable or lower token cost.
The best configuration (\method{}-8) achieves XX.X\% accuracy using an
average of XX.X tokens---a \textbf{+XX.X\,pp} improvement over Fixed-512
at \textbf{XX\%} of the token budget.
Second, \method{} substantially outperforms the compute-matched
self-consistency baseline (SC@$K{\times}128$), confirming that the
consensus signal is more effectively leveraged through adaptive routing
than through simple majority voting.
For instance, at $K = 4$, \method{}-4 achieves XX.X\% vs.\
SC@$4{\times}128$'s XX.X\%, while consuming fewer total tokens because
easy samples are resolved without a resolution phase.
Third, the route distribution reveals that a substantial fraction of
GSM8K problems (XX.X\%) are routed as easy, validating the central
premise that many problems can be solved with minimal compute.

% ---------------------------------------------------------------------------
\subsection{Ablation Studies}
\label{sec:exp:ablation}
% ---------------------------------------------------------------------------

We systematically ablate the three key hyperparameters of \method{} to
understand their individual contributions.

% --- K ablation ---
\paragraph{Effect of parallelism $K$.}
Table~\ref{tab:ablation-k} fixes $B_{\mathrm{probe}} = 128$ and varies
$K \in \{2, 3, 4, 8\}$.
Increasing $K$ improves the reliability of the consensus signal
(Proposition~\ref{prop:monotone}), yielding monotonically higher accuracy.
However, the exploration cost scales linearly in $K$, creating a
diminishing-returns trade-off: from $K = 4$ to $K = 8$, accuracy improves
by XX.X\,pp while average tokens increase by XX.X.
We find $K = 4$ to be the best cost-accuracy operating point for GSM8K.

\begin{table}[t]
\centering
\caption{%
  \textbf{Ablation on $K$} (parallelism) with $B_{\mathrm{probe}} = 128$,
  $B_{\mathrm{med}} = 256$, $B_{\mathrm{hard}} = 512$, default thresholds.
}
\label{tab:ablation-k}
\small
\begin{tabular}{@{}ccccccc@{}}
\toprule
$K$ & \textbf{Acc (\%)} & \textbf{Avg Tok} &
  \textbf{E\%} & \textbf{M\%} & \textbf{H\%} \\
\midrule
2 & XX.X & XX.X & XX.X & XX.X & XX.X \\
3 & XX.X & XX.X & XX.X & XX.X & XX.X \\
4 & XX.X & XX.X & XX.X & XX.X & XX.X \\
8 & XX.X & XX.X & XX.X & XX.X & XX.X \\
\bottomrule
\end{tabular}
\end{table}

% --- Probe budget ablation ---
\paragraph{Effect of probe budget $B_{\mathrm{probe}}$.}
Table~\ref{tab:ablation-probe} fixes $K = 2$ and varies
$B_{\mathrm{probe}} \in \{128, 170, 256, 512\}$.
Larger probe budgets allow the exploration paths to develop more
complete reasoning chains, increasing both the validity rate $\nu$ and
the reliability of the consensus signal.
However, beyond $B_{\mathrm{probe}} = 256$, the additional cost is not
justified: the marginal accuracy gain is XX.X\,pp while the per-sample
cost increases by XX.X tokens.
At $B_{\mathrm{probe}} = 512$, the exploration phase alone consumes as
many tokens as a single Fixed-512 run, eliminating the efficiency
advantage of adaptive routing.

\begin{table}[t]
\centering
\caption{%
  \textbf{Ablation on $B_{\mathrm{probe}}$} with $K = 2$,
  $B_{\mathrm{med}} = 256$, $B_{\mathrm{hard}} = 512$, default thresholds.
}
\label{tab:ablation-probe}
\small
\begin{tabular}{@{}ccccccc@{}}
\toprule
$B_{\mathrm{probe}}$ & \textbf{Acc (\%)} & \textbf{Avg Tok} &
  \textbf{E\%} & \textbf{M\%} & \textbf{H\%} \\
\midrule
128 & XX.X & XX.X & XX.X & XX.X & XX.X \\
170 & XX.X & XX.X & XX.X & XX.X & XX.X \\
256 & XX.X & XX.X & XX.X & XX.X & XX.X \\
512 & XX.X & XX.X & XX.X & XX.X & XX.X \\
\bottomrule
\end{tabular}
\end{table}

% --- Threshold ablation ---
\paragraph{Sensitivity to routing thresholds.}
Table~\ref{tab:ablation-thresh} fixes $K = 4$, $B_{\mathrm{probe}} = 128$
and sweeps the easy-route threshold
$\theta_E \in \{0.75, 0.90, 0.95\}$ and medium-route threshold
$\theta_M \in \{0.50, 0.60, 0.70\}$.
Higher $\theta_E$ makes the router more conservative, sending fewer
problems to the easy route and increasing average cost.
The accuracy is relatively insensitive to $\theta_E$ in the range
$[0.75, 0.95]$, varying by only XX.X\,pp, because problems
mis-classified as easy at $\theta_E = 0.75$ are already near the
decision boundary and contribute minimal accuracy loss.
The medium threshold $\theta_M$ has a larger effect: raising $\theta_M$
from 0.50 to 0.70 shifts samples from the medium to the hard route,
increasing both accuracy (by XX.X\,pp) and cost (by XX.X tokens).

\begin{table}[t]
\centering
\caption{%
  \textbf{Ablation on routing thresholds} with $K = 4$,
  $B_{\mathrm{probe}} = 128$. Each cell: Acc (\%) / Avg Tok.
}
\label{tab:ablation-thresh}
\small
\begin{tabular}{@{}lccc@{}}
\toprule
& \multicolumn{3}{c}{$\theta_M$ (medium threshold)} \\
\cmidrule(lr){2-4}
$\theta_E$ (easy) & 0.50 & 0.60 & 0.70 \\
\midrule
0.75 & XX.X\,/\,XX.X & XX.X\,/\,XX.X & XX.X\,/\,XX.X \\
0.90 & XX.X\,/\,XX.X & XX.X\,/\,XX.X & XX.X\,/\,XX.X \\
0.95 & XX.X\,/\,XX.X & XX.X\,/\,XX.X & XX.X\,/\,XX.X \\
\bottomrule
\end{tabular}
\end{table}

% ---------------------------------------------------------------------------
\subsection{Route Analysis}
\label{sec:exp:routes}
% ---------------------------------------------------------------------------

To understand \emph{how} the consensus router allocates compute, we
analyze the behavior of each difficulty route in detail.

\paragraph{Per-route accuracy.}
Table~\ref{tab:route-accuracy} breaks down accuracy by assigned route
for the best \method{} configuration ($K = 4$, $B_{\mathrm{probe}} = 128$).
Easy-route samples achieve the highest accuracy (XX.X\%), confirming
that the consensus signal reliably identifies problems the model can
solve with minimal compute.
Medium-route accuracy (XX.X\%) is intermediate, while hard-route accuracy
(XX.X\%) is lowest---as expected, since these are the problems where the
model genuinely struggles.
Critically, the easy-route accuracy exceeds what these same samples
achieve under Fixed-512 (XX.X\%), demonstrating that extended reasoning
\emph{hurts} on easy problems (the overthinking phenomenon discussed
in \S\ref{sec:intro}).

\begin{table}[t]
\centering
\caption{%
  \textbf{Per-route analysis} for \method{} ($K = 4$,
  $B_{\mathrm{probe}} = 128$, $n = 200$).
  ``Fixed-512 acc on same'' reports the accuracy of Fixed-512 on the
  subset of samples assigned to each route.
}
\label{tab:route-accuracy}
\small
\begin{tabular}{@{}lcccc@{}}
\toprule
\textbf{Route} & \textbf{\% Samples} & \textbf{Acc (\%)} &
  \textbf{Avg Tok} & \textbf{Fixed-512 acc on same (\%)} \\
\midrule
Easy   & XX.X & XX.X & XX.X & XX.X \\
Medium & XX.X & XX.X & XX.X & XX.X \\
Hard   & XX.X & XX.X & XX.X & XX.X \\
\midrule
Overall & 100.0 & XX.X & XX.X & XX.X \\
\bottomrule
\end{tabular}
\end{table}

\paragraph{Token savings per route.}
The easy route consumes only $K \times B_{\mathrm{probe}}$ tokens
(XX.X on average), saving XX.X\% relative to Fixed-512.
The medium route consumes an average of XX.X tokens (exploration +
extension), while the hard route consumes XX.X tokens (exploration +
full deliberation).
Across all routes, the weighted average token consumption is XX.X,
representing a XX.X\% reduction compared to Fixed-512.

\paragraph{Consensus score and correctness.}
Figure~\ref{fig:consensus-correct} plots the relationship between the
consensus confidence score $\gamma$ and per-sample correctness.
Samples with $\gamma \geq 0.8$ are correct XX.X\% of the time, while
samples with $\gamma < 0.4$ are correct only XX.X\% of the time.
The Spearman rank correlation between $\gamma$ and binary correctness
is $\rho = $ XX.XX ($p < $ XX), empirically validating
Proposition~\ref{prop:monotone}'s prediction that consensus is a
monotone difficulty estimator.
This strong correlation explains why the simple three-tier routing rule
in Eq.~\eqref{eq:routing} is effective: the consensus signal separates
difficulty levels well enough that even coarse binning into three routes
captures most of the available adaptive gain.

\begin{figure}[t]
\centering
% \includegraphics[width=0.55\textwidth]{figures/consensus_vs_correctness.pdf}
\fbox{\parbox{0.55\textwidth}{\centering\vspace{2cm}%
  \textit{Placeholder: scatter/binned plot of consensus score $\gamma$
  vs.\ correctness rate, with fitted logistic curve.}%
\vspace{2cm}}}
\caption{%
  \textbf{Consensus score vs.\ correctness.}
  Each point represents a bin of samples grouped by consensus confidence
  $\gamma$; the $y$-axis shows the fraction of correct answers.
  The strong positive correlation ($\rho =$ XX.XX) confirms that
  multi-path consensus is a reliable difficulty signal.
}
\label{fig:consensus-correct}
\end{figure}


% ---------------------------------------------------------------------------
\subsection{Experimental Setup}
\label{sec:exp:setup}
% ---------------------------------------------------------------------------

\paragraph{Model and inference.}
All experiments use \textbf{Qwen3-8B} with thinking mode enabled
(\texttt{enable\_thinking=True}), which instructs the model to produce an
extended chain-of-thought within \texttt{<think>...</think>} tags before
emitting a final answer.
We control the thinking token budget $b$ by setting
\texttt{max\_thinking\_tokens}; once this limit is reached, the model is
forced to exit the thinking block and produce an answer.
Greedy decoding ($\tau = 0$) is used for the first exploration path and
all resolution phases; the remaining $K{-}1$ exploration paths use
$\tau = 0.7$ to encourage diversity (\S\ref{sec:method:explore}).

\paragraph{Benchmark.}
We evaluate on \textbf{GSM8K}~\citep{cobbe2021training}, a dataset of
1{,}319 grade-school math word problems.
For the main experiments and ablations we sample subsets of $n = 200$
problems (with fixed random seeds for reproducibility); full-dataset
results are reported for the best configuration.
Accuracy is measured via exact match after deterministic number extraction
(the rightmost numerical value in the model's output).

\paragraph{Baselines.}
We compare \method{} against four baselines that span the natural
design space:
\begin{itemize}[nosep,leftmargin=*]
  \item \textbf{Fixed-$b$}: single-pass generation with a fixed thinking
    budget $b \in \{128, 256, 512\}$ tokens.
  \item \textbf{SC@$K{\times}128$}: self-consistency~\citep{wang2023selfconsistency}
    with $K$ independent samples at $b = 128$ tokens each, followed by
    majority voting.
    This baseline uses the same total exploration tokens as \method{}
    with $K$ paths and $B_{\mathrm{probe}} = 128$, providing a
    compute-matched comparison.
\end{itemize}

\paragraph{\method{} configurations.}
We explore the following hyperparameter axes:
parallelism $K \in \{2, 3, 4, 8\}$,
probe budget $B_{\mathrm{probe}} \in \{128, 170, 256, 512\}$,
medium resolution budget $B_{\mathrm{med}} \in \{256, 512\}$, and
hard resolution budget $B_{\mathrm{hard}} \in \{512, 1024\}$.
Routing thresholds default to $\theta_E = 0.75$ (easy) and
$\theta_M = 0.5$ (medium) unless varied in ablation studies.
The consensus confidence weights are fixed at
$\gamma = 0.5\alpha + 0.3\nu + 0.2\mu$ throughout.

\paragraph{Evaluation metrics.}
For each method we report:
(1)~\textbf{Accuracy} (exact-match \%),
(2)~\textbf{Avg Tokens} (mean thinking tokens consumed per sample,
including exploration overhead for \method{}),
(3)~\textbf{Route distribution} (fraction of samples classified as
Easy / Medium / Hard by the consensus router).

\paragraph{Hardware.}
All experiments are conducted on a single NVIDIA A100-80GB GPU.

% ---------------------------------------------------------------------------
\subsection{Main Results}
\label{sec:exp:main}
% ---------------------------------------------------------------------------

Table~\ref{tab:main-rs} presents the primary comparison between
\method{}, fixed-budget baselines, and self-consistency on GSM8K
($n = 200$).
All \method{} variants use default routing thresholds
($\theta_E = 0.75$, $\theta_M = 0.5$).

\begin{table}[t]
\centering
\caption{%
  \textbf{Main results on GSM8K} ($n = 200$, Qwen3-8B thinking mode).
  \textbf{Avg Tok} includes full end-to-end cost (exploration +
  resolution).
  Route columns report the fraction of samples classified as Easy (E),
  Medium (M), and Hard (H) by the consensus router.
  Best adaptive result in \textbf{bold}; best baseline \underline{underlined}.
}
\label{tab:main-rs}
\small
\setlength{\tabcolsep}{4pt}
\begin{tabular}{@{}lcccccccc@{}}
\toprule
\textbf{Method} & $K$ & $B_{\mathrm{probe}}$ &
  \textbf{Acc (\%)} & \textbf{Avg Tok} &
  \textbf{E\%} & \textbf{M\%} & \textbf{H\%} \\
\midrule
\multicolumn{8}{@{}l}{\textit{Fixed-budget baselines}} \\
Fixed-128       & -- & -- & XX.X & 128  & -- & -- & -- \\
Fixed-256       & -- & -- & XX.X & 256  & -- & -- & -- \\
\underline{Fixed-512} & -- & -- & \underline{XX.X} & 512  & -- & -- & -- \\
\midrule
\multicolumn{8}{@{}l}{\textit{Self-consistency baselines}} \\
SC@$2{\times}128$  & 2 & 128 & XX.X & 256  & -- & -- & -- \\
SC@$4{\times}128$  & 4 & 128 & XX.X & 512  & -- & -- & -- \\
SC@$8{\times}128$  & 8 & 128 & XX.X & 1024 & -- & -- & -- \\
\midrule
\multicolumn{8}{@{}l}{\textit{\method{} (ours)}} \\
\method{}-2      & 2 & 128 & XX.X & XX.X & XX.X & XX.X & XX.X \\
\method{}-4      & 4 & 128 & XX.X & XX.X & XX.X & XX.X & XX.X \\
\method{}-4-256  & 4 & 256 & XX.X & XX.X & XX.X & XX.X & XX.X \\
\textbf{\method{}-8} & 8 & 128 & \textbf{XX.X} & XX.X & XX.X & XX.X & XX.X \\
\method{}-8-256  & 8 & 256 & XX.X & XX.X & XX.X & XX.X & XX.X \\
\bottomrule
\end{tabular}
\end{table}

Several observations emerge from Table~\ref{tab:main-rs}.
First, \method{} consistently outperforms fixed-budget baselines at
comparable or lower token cost.
The best configuration (\method{}-8) achieves XX.X\% accuracy using an
average of XX.X tokens---a \textbf{+XX.X\,pp} improvement over Fixed-512
at \textbf{XX\%} of the token budget.
Second, \method{} substantially outperforms the compute-matched
self-consistency baseline (SC@$K{\times}128$), confirming that the
consensus signal is more effectively leveraged through adaptive routing
than through simple majority voting.
For instance, at $K = 4$, \method{}-4 achieves XX.X\% vs.\
SC@$4{\times}128$'s XX.X\%, while consuming fewer total tokens because
easy samples are resolved without a resolution phase.
Third, the route distribution reveals that a substantial fraction of
GSM8K problems (XX.X\%) are routed as easy, validating the central
premise that many problems can be solved with minimal compute.

% ---------------------------------------------------------------------------
\subsection{Ablation Studies}
\label{sec:exp:ablation}
% ---------------------------------------------------------------------------

We systematically ablate the three key hyperparameters of \method{} to
understand their individual contributions.

% --- K ablation ---
\paragraph{Effect of parallelism $K$.}
Table~\ref{tab:ablation-k} fixes $B_{\mathrm{probe}} = 128$ and varies
$K \in \{2, 3, 4, 8\}$.
Increasing $K$ improves the reliability of the consensus signal
(Proposition~\ref{prop:monotone}), yielding monotonically higher accuracy.
However, the exploration cost scales linearly in $K$, creating a
diminishing-returns trade-off: from $K = 4$ to $K = 8$, accuracy improves
by XX.X\,pp while average tokens increase by XX.X.
We find $K = 4$ to be the best cost-accuracy operating point for GSM8K.

\begin{table}[t]
\centering
\caption{%
  \textbf{Ablation on $K$} (parallelism) with $B_{\mathrm{probe}} = 128$,
  $B_{\mathrm{med}} = 256$, $B_{\mathrm{hard}} = 512$, default thresholds.
}
\label{tab:ablation-k}
\small
\begin{tabular}{@{}ccccccc@{}}
\toprule
$K$ & \textbf{Acc (\%)} & \textbf{Avg Tok} &
  \textbf{E\%} & \textbf{M\%} & \textbf{H\%} \\
\midrule
2 & XX.X & XX.X & XX.X & XX.X & XX.X \\
3 & XX.X & XX.X & XX.X & XX.X & XX.X \\
4 & XX.X & XX.X & XX.X & XX.X & XX.X \\
8 & XX.X & XX.X & XX.X & XX.X & XX.X \\
\bottomrule
\end{tabular}
\end{table}

% --- Probe budget ablation ---
\paragraph{Effect of probe budget $B_{\mathrm{probe}}$.}
Table~\ref{tab:ablation-probe} fixes $K = 2$ and varies
$B_{\mathrm{probe}} \in \{128, 170, 256, 512\}$.
Larger probe budgets allow the exploration paths to develop more
complete reasoning chains, increasing both the validity rate $\nu$ and
the reliability of the consensus signal.
However, beyond $B_{\mathrm{probe}} = 256$, the additional cost is not
justified: the marginal accuracy gain is XX.X\,pp while the per-sample
cost increases by XX.X tokens.
At $B_{\mathrm{probe}} = 512$, the exploration phase alone consumes as
many tokens as a single Fixed-512 run, eliminating the efficiency
advantage of adaptive routing.

\begin{table}[t]
\centering
\caption{%
  \textbf{Ablation on $B_{\mathrm{probe}}$} with $K = 2$,
  $B_{\mathrm{med}} = 256$, $B_{\mathrm{hard}} = 512$, default thresholds.
}
\label{tab:ablation-probe}
\small
\begin{tabular}{@{}ccccccc@{}}
\toprule
$B_{\mathrm{probe}}$ & \textbf{Acc (\%)} & \textbf{Avg Tok} &
  \textbf{E\%} & \textbf{M\%} & \textbf{H\%} \\
\midrule
128 & XX.X & XX.X & XX.X & XX.X & XX.X \\
170 & XX.X & XX.X & XX.X & XX.X & XX.X \\
256 & XX.X & XX.X & XX.X & XX.X & XX.X \\
512 & XX.X & XX.X & XX.X & XX.X & XX.X \\
\bottomrule
\end{tabular}
\end{table}

% --- Threshold ablation ---
\paragraph{Sensitivity to routing thresholds.}
Table~\ref{tab:ablation-thresh} fixes $K = 4$, $B_{\mathrm{probe}} = 128$
and sweeps the easy-route threshold
$\theta_E \in \{0.75, 0.90, 0.95\}$ and medium-route threshold
$\theta_M \in \{0.50, 0.60, 0.70\}$.
Higher $\theta_E$ makes the router more conservative, sending fewer
problems to the easy route and increasing average cost.
The accuracy is relatively insensitive to $\theta_E$ in the range
$[0.75, 0.95]$, varying by only XX.X\,pp, because problems
mis-classified as easy at $\theta_E = 0.75$ are already near the
decision boundary and contribute minimal accuracy loss.
The medium threshold $\theta_M$ has a larger effect: raising $\theta_M$
from 0.50 to 0.70 shifts samples from the medium to the hard route,
increasing both accuracy (by XX.X\,pp) and cost (by XX.X tokens).

\begin{table}[t]
\centering
\caption{%
  \textbf{Ablation on routing thresholds} with $K = 4$,
  $B_{\mathrm{probe}} = 128$. Each cell: Acc (\%) / Avg Tok.
}
\label{tab:ablation-thresh}
\small
\begin{tabular}{@{}lccc@{}}
\toprule
& \multicolumn{3}{c}{$\theta_M$ (medium threshold)} \\
\cmidrule(lr){2-4}
$\theta_E$ (easy) & 0.50 & 0.60 & 0.70 \\
\midrule
0.75 & XX.X\,/\,XX.X & XX.X\,/\,XX.X & XX.X\,/\,XX.X \\
0.90 & XX.X\,/\,XX.X & XX.X\,/\,XX.X & XX.X\,/\,XX.X \\
0.95 & XX.X\,/\,XX.X & XX.X\,/\,XX.X & XX.X\,/\,XX.X \\
\bottomrule
\end{tabular}
\end{table}

% ---------------------------------------------------------------------------
\subsection{Route Analysis}
\label{sec:exp:routes}
% ---------------------------------------------------------------------------

To understand \emph{how} the consensus router allocates compute, we
analyze the behavior of each difficulty route in detail.

\paragraph{Per-route accuracy.}
Table~\ref{tab:route-accuracy} breaks down accuracy by assigned route
for the best \method{} configuration ($K = 4$, $B_{\mathrm{probe}} = 128$).
Easy-route samples achieve the highest accuracy (XX.X\%), confirming
that the consensus signal reliably identifies problems the model can
solve with minimal compute.
Medium-route accuracy (XX.X\%) is intermediate, while hard-route accuracy
(XX.X\%) is lowest---as expected, since these are the problems where the
model genuinely struggles.
Critically, the easy-route accuracy exceeds what these same samples
achieve under Fixed-512 (XX.X\%), demonstrating that extended reasoning
\emph{hurts} on easy problems (the overthinking phenomenon discussed
in \S\ref{sec:intro}).

\begin{table}[t]
\centering
\caption{%
  \textbf{Per-route analysis} for \method{} ($K = 4$,
  $B_{\mathrm{probe}} = 128$, $n = 200$).
  ``Fixed-512 acc on same'' reports the accuracy of Fixed-512 on the
  subset of samples assigned to each route.
}
\label{tab:route-accuracy}
\small
\begin{tabular}{@{}lcccc@{}}
\toprule
\textbf{Route} & \textbf{\% Samples} & \textbf{Acc (\%)} &
  \textbf{Avg Tok} & \textbf{Fixed-512 acc on same (\%)} \\
\midrule
Easy   & XX.X & XX.X & XX.X & XX.X \\
Medium & XX.X & XX.X & XX.X & XX.X \\
Hard   & XX.X & XX.X & XX.X & XX.X \\
\midrule
Overall & 100.0 & XX.X & XX.X & XX.X \\
\bottomrule
\end{tabular}
\end{table}

\paragraph{Token savings per route.}
The easy route consumes only $K \times B_{\mathrm{probe}}$ tokens
(XX.X on average), saving XX.X\% relative to Fixed-512.
The medium route consumes an average of XX.X tokens (exploration +
extension), while the hard route consumes XX.X tokens (exploration +
full deliberation).
Across all routes, the weighted average token consumption is XX.X,
representing a XX.X\% reduction compared to Fixed-512.

\paragraph{Consensus score and correctness.}
Figure~\ref{fig:consensus-correct} plots the relationship between the
consensus confidence score $\gamma$ and per-sample correctness.
Samples with $\gamma \geq 0.8$ are correct XX.X\% of the time, while
samples with $\gamma < 0.4$ are correct only XX.X\% of the time.
The Spearman rank correlation between $\gamma$ and binary correctness
is $\rho = $ XX.XX ($p < $ XX), empirically validating
Proposition~\ref{prop:monotone}'s prediction that consensus is a
monotone difficulty estimator.
This strong correlation explains why the simple three-tier routing rule
in Eq.~\eqref{eq:routing} is effective: the consensus signal separates
difficulty levels well enough that even coarse binning into three routes
captures most of the available adaptive gain.

\begin{figure}[t]
\centering
% \includegraphics[width=0.55\textwidth]{figures/consensus_vs_correctness.pdf}
\fbox{\parbox{0.55\textwidth}{\centering\vspace{2cm}%
  \textit{Placeholder: scatter/binned plot of consensus score $\gamma$
  vs.\ correctness rate, with fitted logistic curve.}%
\vspace{2cm}}}
\caption{%
  \textbf{Consensus score vs.\ correctness.}
  Each point represents a bin of samples grouped by consensus confidence
  $\gamma$; the $y$-axis shows the fraction of correct answers.
  The strong positive correlation ($\rho =$ XX.XX) confirms that
  multi-path consensus is a reliable difficulty signal.
}
\label{fig:consensus-correct}
\end{figure}


% ---------------------------------------------------------------------------
\subsection{Experimental Setup}
\label{sec:exp:setup}
% ---------------------------------------------------------------------------

\paragraph{Model and inference.}
All experiments use \textbf{Qwen3-8B} with thinking mode enabled
(\texttt{enable\_thinking=True}), which instructs the model to produce an
extended chain-of-thought within \texttt{<think>...</think>} tags before
emitting a final answer.
We control the thinking token budget $b$ by setting
\texttt{max\_thinking\_tokens}; once this limit is reached, the model is
forced to exit the thinking block and produce an answer.
Greedy decoding ($\tau = 0$) is used for the first exploration path and
all resolution phases; the remaining $K{-}1$ exploration paths use
$\tau = 0.7$ to encourage diversity (\S\ref{sec:method:explore}).

\paragraph{Benchmark.}
We evaluate on \textbf{GSM8K}~\citep{cobbe2021training}, a dataset of
1{,}319 grade-school math word problems.
For the main experiments and ablations we sample subsets of $n = 200$
problems (with fixed random seeds for reproducibility); full-dataset
results are reported for the best configuration.
Accuracy is measured via exact match after deterministic number extraction
(the rightmost numerical value in the model's output).

\paragraph{Baselines.}
We compare \method{} against four baselines that span the natural
design space:
\begin{itemize}[nosep,leftmargin=*]
  \item \textbf{Fixed-$b$}: single-pass generation with a fixed thinking
    budget $b \in \{128, 256, 512\}$ tokens.
  \item \textbf{SC@$K{\times}128$}: self-consistency~\citep{wang2023selfconsistency}
    with $K$ independent samples at $b = 128$ tokens each, followed by
    majority voting.
    This baseline uses the same total exploration tokens as \method{}
    with $K$ paths and $B_{\mathrm{probe}} = 128$, providing a
    compute-matched comparison.
\end{itemize}

\paragraph{\method{} configurations.}
We explore the following hyperparameter axes:
parallelism $K \in \{2, 3, 4, 8\}$,
probe budget $B_{\mathrm{probe}} \in \{128, 170, 256, 512\}$,
medium resolution budget $B_{\mathrm{med}} \in \{256, 512\}$, and
hard resolution budget $B_{\mathrm{hard}} \in \{512, 1024\}$.
Routing thresholds default to $\theta_E = 0.75$ (easy) and
$\theta_M = 0.5$ (medium) unless varied in ablation studies.
The consensus confidence weights are fixed at
$\gamma = 0.5\alpha + 0.3\nu + 0.2\mu$ throughout.

\paragraph{Evaluation metrics.}
For each method we report:
(1)~\textbf{Accuracy} (exact-match \%),
(2)~\textbf{Avg Tokens} (mean thinking tokens consumed per sample,
including exploration overhead for \method{}),
(3)~\textbf{Route distribution} (fraction of samples classified as
Easy / Medium / Hard by the consensus router).

\paragraph{Hardware.}
All experiments are conducted on a single NVIDIA A100-80GB GPU.

% ---------------------------------------------------------------------------
\subsection{Main Results}
\label{sec:exp:main}
% ---------------------------------------------------------------------------

Table~\ref{tab:main-rs} presents the primary comparison between
\method{}, fixed-budget baselines, and self-consistency on GSM8K
($n = 200$).
All \method{} variants use default routing thresholds
($\theta_E = 0.75$, $\theta_M = 0.5$).

\begin{table}[t]
\centering
\caption{%
  \textbf{Main results on GSM8K} ($n = 200$, Qwen3-8B thinking mode).
  \textbf{Avg Tok} includes full end-to-end cost (exploration +
  resolution).
  Route columns report the fraction of samples classified as Easy (E),
  Medium (M), and Hard (H) by the consensus router.
  Best adaptive result in \textbf{bold}; best baseline \underline{underlined}.
}
\label{tab:main-rs}
\small
\setlength{\tabcolsep}{4pt}
\begin{tabular}{@{}lcccccccc@{}}
\toprule
\textbf{Method} & $K$ & $B_{\mathrm{probe}}$ &
  \textbf{Acc (\%)} & \textbf{Avg Tok} &
  \textbf{E\%} & \textbf{M\%} & \textbf{H\%} \\
\midrule
\multicolumn{8}{@{}l}{\textit{Fixed-budget baselines}} \\
Fixed-128       & -- & -- & XX.X & 128  & -- & -- & -- \\
Fixed-256       & -- & -- & XX.X & 256  & -- & -- & -- \\
\underline{Fixed-512} & -- & -- & \underline{XX.X} & 512  & -- & -- & -- \\
\midrule
\multicolumn{8}{@{}l}{\textit{Self-consistency baselines}} \\
SC@$2{\times}128$  & 2 & 128 & XX.X & 256  & -- & -- & -- \\
SC@$4{\times}128$  & 4 & 128 & XX.X & 512  & -- & -- & -- \\
SC@$8{\times}128$  & 8 & 128 & XX.X & 1024 & -- & -- & -- \\
\midrule
\multicolumn{8}{@{}l}{\textit{\method{} (ours)}} \\
\method{}-2      & 2 & 128 & XX.X & XX.X & XX.X & XX.X & XX.X \\
\method{}-4      & 4 & 128 & XX.X & XX.X & XX.X & XX.X & XX.X \\
\method{}-4-256  & 4 & 256 & XX.X & XX.X & XX.X & XX.X & XX.X \\
\textbf{\method{}-8} & 8 & 128 & \textbf{XX.X} & XX.X & XX.X & XX.X & XX.X \\
\method{}-8-256  & 8 & 256 & XX.X & XX.X & XX.X & XX.X & XX.X \\
\bottomrule
\end{tabular}
\end{table}

Several observations emerge from Table~\ref{tab:main-rs}.
First, \method{} consistently outperforms fixed-budget baselines at
comparable or lower token cost.
The best configuration (\method{}-8) achieves XX.X\% accuracy using an
average of XX.X tokens---a \textbf{+XX.X\,pp} improvement over Fixed-512
at \textbf{XX\%} of the token budget.
Second, \method{} substantially outperforms the compute-matched
self-consistency baseline (SC@$K{\times}128$), confirming that the
consensus signal is more effectively leveraged through adaptive routing
than through simple majority voting.
For instance, at $K = 4$, \method{}-4 achieves XX.X\% vs.\
SC@$4{\times}128$'s XX.X\%, while consuming fewer total tokens because
easy samples are resolved without a resolution phase.
Third, the route distribution reveals that a substantial fraction of
GSM8K problems (XX.X\%) are routed as easy, validating the central
premise that many problems can be solved with minimal compute.

% ---------------------------------------------------------------------------
\subsection{Ablation Studies}
\label{sec:exp:ablation}
% ---------------------------------------------------------------------------

We systematically ablate the three key hyperparameters of \method{} to
understand their individual contributions.

% --- K ablation ---
\paragraph{Effect of parallelism $K$.}
Table~\ref{tab:ablation-k} fixes $B_{\mathrm{probe}} = 128$ and varies
$K \in \{2, 3, 4, 8\}$.
Increasing $K$ improves the reliability of the consensus signal
(Proposition~\ref{prop:monotone}), yielding monotonically higher accuracy.
However, the exploration cost scales linearly in $K$, creating a
diminishing-returns trade-off: from $K = 4$ to $K = 8$, accuracy improves
by XX.X\,pp while average tokens increase by XX.X.
We find $K = 4$ to be the best cost-accuracy operating point for GSM8K.

\begin{table}[t]
\centering
\caption{%
  \textbf{Ablation on $K$} (parallelism) with $B_{\mathrm{probe}} = 128$,
  $B_{\mathrm{med}} = 256$, $B_{\mathrm{hard}} = 512$, default thresholds.
}
\label{tab:ablation-k}
\small
\begin{tabular}{@{}ccccccc@{}}
\toprule
$K$ & \textbf{Acc (\%)} & \textbf{Avg Tok} &
  \textbf{E\%} & \textbf{M\%} & \textbf{H\%} \\
\midrule
2 & XX.X & XX.X & XX.X & XX.X & XX.X \\
3 & XX.X & XX.X & XX.X & XX.X & XX.X \\
4 & XX.X & XX.X & XX.X & XX.X & XX.X \\
8 & XX.X & XX.X & XX.X & XX.X & XX.X \\
\bottomrule
\end{tabular}
\end{table}

% --- Probe budget ablation ---
\paragraph{Effect of probe budget $B_{\mathrm{probe}}$.}
Table~\ref{tab:ablation-probe} fixes $K = 2$ and varies
$B_{\mathrm{probe}} \in \{128, 170, 256, 512\}$.
Larger probe budgets allow the exploration paths to develop more
complete reasoning chains, increasing both the validity rate $\nu$ and
the reliability of the consensus signal.
However, beyond $B_{\mathrm{probe}} = 256$, the additional cost is not
justified: the marginal accuracy gain is XX.X\,pp while the per-sample
cost increases by XX.X tokens.
At $B_{\mathrm{probe}} = 512$, the exploration phase alone consumes as
many tokens as a single Fixed-512 run, eliminating the efficiency
advantage of adaptive routing.

\begin{table}[t]
\centering
\caption{%
  \textbf{Ablation on $B_{\mathrm{probe}}$} with $K = 2$,
  $B_{\mathrm{med}} = 256$, $B_{\mathrm{hard}} = 512$, default thresholds.
}
\label{tab:ablation-probe}
\small
\begin{tabular}{@{}ccccccc@{}}
\toprule
$B_{\mathrm{probe}}$ & \textbf{Acc (\%)} & \textbf{Avg Tok} &
  \textbf{E\%} & \textbf{M\%} & \textbf{H\%} \\
\midrule
128 & XX.X & XX.X & XX.X & XX.X & XX.X \\
170 & XX.X & XX.X & XX.X & XX.X & XX.X \\
256 & XX.X & XX.X & XX.X & XX.X & XX.X \\
512 & XX.X & XX.X & XX.X & XX.X & XX.X \\
\bottomrule
\end{tabular}
\end{table}

% --- Threshold ablation ---
\paragraph{Sensitivity to routing thresholds.}
Table~\ref{tab:ablation-thresh} fixes $K = 4$, $B_{\mathrm{probe}} = 128$
and sweeps the easy-route threshold
$\theta_E \in \{0.75, 0.90, 0.95\}$ and medium-route threshold
$\theta_M \in \{0.50, 0.60, 0.70\}$.
Higher $\theta_E$ makes the router more conservative, sending fewer
problems to the easy route and increasing average cost.
The accuracy is relatively insensitive to $\theta_E$ in the range
$[0.75, 0.95]$, varying by only XX.X\,pp, because problems
mis-classified as easy at $\theta_E = 0.75$ are already near the
decision boundary and contribute minimal accuracy loss.
The medium threshold $\theta_M$ has a larger effect: raising $\theta_M$
from 0.50 to 0.70 shifts samples from the medium to the hard route,
increasing both accuracy (by XX.X\,pp) and cost (by XX.X tokens).

\begin{table}[t]
\centering
\caption{%
  \textbf{Ablation on routing thresholds} with $K = 4$,
  $B_{\mathrm{probe}} = 128$. Each cell: Acc (\%) / Avg Tok.
}
\label{tab:ablation-thresh}
\small
\begin{tabular}{@{}lccc@{}}
\toprule
& \multicolumn{3}{c}{$\theta_M$ (medium threshold)} \\
\cmidrule(lr){2-4}
$\theta_E$ (easy) & 0.50 & 0.60 & 0.70 \\
\midrule
0.75 & XX.X\,/\,XX.X & XX.X\,/\,XX.X & XX.X\,/\,XX.X \\
0.90 & XX.X\,/\,XX.X & XX.X\,/\,XX.X & XX.X\,/\,XX.X \\
0.95 & XX.X\,/\,XX.X & XX.X\,/\,XX.X & XX.X\,/\,XX.X \\
\bottomrule
\end{tabular}
\end{table}

% ---------------------------------------------------------------------------
\subsection{Route Analysis}
\label{sec:exp:routes}
% ---------------------------------------------------------------------------

To understand \emph{how} the consensus router allocates compute, we
analyze the behavior of each difficulty route in detail.

\paragraph{Per-route accuracy.}
Table~\ref{tab:route-accuracy} breaks down accuracy by assigned route
for the best \method{} configuration ($K = 4$, $B_{\mathrm{probe}} = 128$).
Easy-route samples achieve the highest accuracy (XX.X\%), confirming
that the consensus signal reliably identifies problems the model can
solve with minimal compute.
Medium-route accuracy (XX.X\%) is intermediate, while hard-route accuracy
(XX.X\%) is lowest---as expected, since these are the problems where the
model genuinely struggles.
Critically, the easy-route accuracy exceeds what these same samples
achieve under Fixed-512 (XX.X\%), demonstrating that extended reasoning
\emph{hurts} on easy problems (the overthinking phenomenon discussed
in \S\ref{sec:intro}).

\begin{table}[t]
\centering
\caption{%
  \textbf{Per-route analysis} for \method{} ($K = 4$,
  $B_{\mathrm{probe}} = 128$, $n = 200$).
  ``Fixed-512 acc on same'' reports the accuracy of Fixed-512 on the
  subset of samples assigned to each route.
}
\label{tab:route-accuracy}
\small
\begin{tabular}{@{}lcccc@{}}
\toprule
\textbf{Route} & \textbf{\% Samples} & \textbf{Acc (\%)} &
  \textbf{Avg Tok} & \textbf{Fixed-512 acc on same (\%)} \\
\midrule
Easy   & XX.X & XX.X & XX.X & XX.X \\
Medium & XX.X & XX.X & XX.X & XX.X \\
Hard   & XX.X & XX.X & XX.X & XX.X \\
\midrule
Overall & 100.0 & XX.X & XX.X & XX.X \\
\bottomrule
\end{tabular}
\end{table}

\paragraph{Token savings per route.}
The easy route consumes only $K \times B_{\mathrm{probe}}$ tokens
(XX.X on average), saving XX.X\% relative to Fixed-512.
The medium route consumes an average of XX.X tokens (exploration +
extension), while the hard route consumes XX.X tokens (exploration +
full deliberation).
Across all routes, the weighted average token consumption is XX.X,
representing a XX.X\% reduction compared to Fixed-512.

\paragraph{Consensus score and correctness.}
Figure~\ref{fig:consensus-correct} plots the relationship between the
consensus confidence score $\gamma$ and per-sample correctness.
Samples with $\gamma \geq 0.8$ are correct XX.X\% of the time, while
samples with $\gamma < 0.4$ are correct only XX.X\% of the time.
The Spearman rank correlation between $\gamma$ and binary correctness
is $\rho = $ XX.XX ($p < $ XX), empirically validating
Proposition~\ref{prop:monotone}'s prediction that consensus is a
monotone difficulty estimator.
This strong correlation explains why the simple three-tier routing rule
in Eq.~\eqref{eq:routing} is effective: the consensus signal separates
difficulty levels well enough that even coarse binning into three routes
captures most of the available adaptive gain.

\begin{figure}[t]
\centering
% \includegraphics[width=0.55\textwidth]{figures/consensus_vs_correctness.pdf}
\fbox{\parbox{0.55\textwidth}{\centering\vspace{2cm}%
  \textit{Placeholder: scatter/binned plot of consensus score $\gamma$
  vs.\ correctness rate, with fitted logistic curve.}%
\vspace{2cm}}}
\caption{%
  \textbf{Consensus score vs.\ correctness.}
  Each point represents a bin of samples grouped by consensus confidence
  $\gamma$; the $y$-axis shows the fraction of correct answers.
  The strong positive correlation ($\rho =$ XX.XX) confirms that
  multi-path consensus is a reliable difficulty signal.
}
\label{fig:consensus-correct}
\end{figure}


% ============================================================================
\section{Analysis}
\label{sec:analysis}
% ============================================================================

% Analysis section — deeper dives beyond main results

\subsection{Why Do Single-Path Controllers Fail?}
\label{sec:analysis:single_path}

A natural question is why simpler approaches---extracting difficulty signals from a single reasoning trace---consistently fail.
Table~\ref{tab:negative_results} summarizes our systematic evaluation of single-path controllers.

\begin{table}[h]
\centering
\caption{Systematic negative results for single-path difficulty controllers on GSM8K.
All methods use the same model (Qwen3.5-27B) and evaluation protocol.
$\Delta$ is accuracy change relative to the best fixed baseline.}
\label{tab:negative_results}
\begin{tabular}{lcc}
\toprule
Controller & Accuracy (\%) & $\Delta$ (pp) \\
\midrule
Best Fixed Budget (256 tokens) & XX.X & --- \\
\midrule
\textit{Feature-based:} \\
\quad Answer Presence + Token Util. + Consistency & XX.X & $-6.5$ \\
\quad Token Utilization Only & XX.X & $-4.2$ \\
\quad Answer Presence Only & XX.X & $-3.1$ \\
\midrule
\textit{Uncertainty-based:} \\
\quad Logit Entropy (requires white-box) & XX.X & $-6.4$ \\
\midrule
\textit{Search-based:} \\
\quad MCTS-R (serial explore-refine-verify) & XX.X & $-5.6$ \\
\midrule
\textbf{\method{} (ours, K=XX, probe=XXX)} & \textbf{XX.X} & \textbf{+X.X} \\
\bottomrule
\end{tabular}
\end{table}

The fundamental issue is \emph{information poverty}: a single reasoning trace provides at most $H(X_1)$ bits of information about problem difficulty, while $K$ independent traces provide up to $K \cdot H(X_1)$ bits (Proposition~\ref{prop:information}).
In practice, the information gain from single-path features is severely limited because:
(i)~answer presence is a noisy binary signal confounded by generation length;
(ii)~token utilization conflates difficulty with verbosity;
(iii)~logit-based uncertainty requires white-box access unavailable in API settings.

\subsection{Consensus Signal Quality}
\label{sec:analysis:consensus}

We validate the consensus signal against ground-truth difficulty using our Phase~0 analysis on existing multi-budget data (Section~\ref{sec:theory}).

\paragraph{Correlation with difficulty.}
The Spearman rank correlation between consensus score $C_K$ and problem difficulty is $\rho = -0.51$ ($p < 10^{-88}$), substantially stronger than any single-path feature (all $|\rho| < 0.2$).
This confirms Proposition~\ref{prop:information}: multi-path consensus is a fundamentally richer difficulty signal.

\paragraph{Consensus precision.}
When all $K$ paths agree on the same answer (full consensus), the answer is correct \textbf{72.8\%} of the time.
This high precision justifies the ``easy'' route: accepting the consensus answer without further computation.

\paragraph{Route calibration.}
Figure~\ref{fig:route_calibration} shows the accuracy of each route as a function of the consensus score $C_K$.
The easy route (high consensus) achieves XX.X\% accuracy, the medium route achieves XX.X\%, and the hard route achieves XX.X\%.
The monotonic relationship confirms that our threshold-based routing is well-calibrated.

\begin{figure}[h]
\centering
% PLACEHOLDER — will be generated from results
\fbox{\parbox{0.8\textwidth}{\centering\vspace{2cm}\textit{[Figure: Route calibration — accuracy vs consensus score, colored by route]}\vspace{2cm}}}
\caption{Accuracy as a function of consensus score. Higher consensus reliably indicates easier problems, validating the routing mechanism.}
\label{fig:route_calibration}
\end{figure}

\subsection{Token Efficiency Analysis}
\label{sec:analysis:efficiency}

\method{} achieves token savings primarily through the easy route, which avoids the costly resolution phase entirely.

\paragraph{Per-route token cost.}
\begin{itemize}[nosep]
\item \textbf{Easy} (XX\% of samples): $K \times B_\text{probe} = $ XX tokens. \textit{No resolution cost.}
\item \textbf{Medium} (XX\% of samples): $K \times B_\text{probe} + B_\text{ext} \approx $ XX tokens.
\item \textbf{Hard} (XX\% of samples): $K \times B_\text{probe} + B_\text{hard} = $ XX tokens.
\end{itemize}

\paragraph{Comparison to budget-matched baselines.}
A fair comparison accounts for the total tokens consumed.
\method{} uses an average of XX tokens per sample.
A fixed budget of XX tokens (the same average) achieves only XX.X\% accuracy, while \method{} achieves XX.X\%---a \textbf{+X.X\,pp} improvement at the same token budget.
This demonstrates that \method{}'s gains come from \emph{better allocation}, not simply \emph{more tokens}.

\paragraph{Wall-clock latency.}
Although \method{} generates $K$ exploration paths, these are \emph{embarrassingly parallel}: on hardware with sufficient memory bandwidth, the $K$ paths can be generated as a single batch, reducing wall-clock latency to approximately $B_\text{probe}$ tokens (not $K \times B_\text{probe}$).
In our sequential implementation, the latency overhead is $K\times$ for the exploration phase, but the resolution phase uses only a single generation call.

\subsection{Error Analysis}
\label{sec:analysis:errors}

We categorize errors by route to understand where \method{} fails:

\begin{itemize}[nosep]
\item \textbf{Easy-route errors} (false confidence): All paths agree on the \emph{wrong} answer.
This occurs when the problem has a common misconception that tricks the model consistently.
These errors are \emph{inherent}---no amount of additional compute would help, as the model lacks the capability.

\item \textbf{Medium-route errors} (insufficient extension): The extension phase fails to correct the initial partial answer.
Often caused by the extension continuing a flawed reasoning chain rather than starting fresh.

\item \textbf{Hard-route errors} (deliberation failure): Despite seeing all path summaries, the deliberation produces an incorrect answer.
These represent genuinely difficult problems where the model struggles even with full context.
\end{itemize}


% ============================================================================
\section{Conclusion}
\label{sec:conclusion}
% ============================================================================

% conclusion_reasonspec.tex — Conclusion for the Reasoning Speculation paper
% Included from main_reasonspec.tex via % conclusion_reasonspec.tex — Conclusion for the Reasoning Speculation paper
% Included from main_reasonspec.tex via % conclusion_reasonspec.tex — Conclusion for the Reasoning Speculation paper
% Included from main_reasonspec.tex via \input{sections/conclusion_reasonspec}

We introduced \textbf{\fullmethod{}} (\method{}), a training-free framework that restructures test-time reasoning compute via parallel exploration and consensus-based routing.
Rather than allocating a single token budget and relying on weak single-path signals to judge difficulty, \method{} generates $K$ short exploration paths, computes a cross-path consensus score, and adaptively routes each problem through one of three resolution strategies---accept, extend, or deliberate.
The key insight underlying our approach is that multi-path consensus is a \emph{fundamentally richer} difficulty signal than any feature extractable from a single reasoning trace: it captures inter-path variation that is information-theoretically inaccessible to single-path methods.

\paragraph{Key findings.}
Our work yields both negative and positive results of independent interest.
On the negative side, we conducted a comprehensive evaluation of all major categories of single-path adaptive budget controllers---honest features, uncertainty-based routing, and MCTS-style refinement---and found that \emph{every} controller degrades accuracy by 5--7\,pp relative to a fixed budget, demonstrating a fundamental limitation of single-path difficulty estimation in the black-box setting.
These systematic negative results motivate the shift to multi-path computation.
On the positive side, we showed that the consensus signal achieves a Spearman correlation of $\rho = -0.51$ with oracle difficulty ($p < 10^{-88}$), far exceeding any single-path feature ($|\rho| \le 0.19$).
We proved that the optimal budget-allocation policy under standard regularity conditions is a threshold policy on consensus, and that the three-route design of \method{} naturally emerges from this optimality structure (Proposition~\ref{prop:threshold}).

\paragraph{Limitations.}
Several limitations of the current work deserve acknowledgment.
First, even though the $K$ exploration paths are embarrassingly parallel and wall-clock latency remains comparable to a single full-budget pass, the \emph{total token count} increases by a factor of up to $K$ for easy questions---a cost that matters when billing is per-token rather than per-second.
Second, the consensus mechanism currently operates only on final extracted answers and does not exploit agreement or divergence in intermediate reasoning steps, potentially discarding useful signal.
Third, our evaluation focuses primarily on mathematical reasoning (GSM8K); generalization to other reasoning domains---code generation, commonsense reasoning, scientific problem-solving---remains to be validated.
Finally, the hyperparameters $K$, $\theta_{\mathrm{E}}$, and $\theta_{\mathrm{M}}$ were tuned on a held-out validation split of GSM8K, and may require task-specific adjustment for new domains.

\paragraph{Future directions.}
The \method{} framework opens several avenues for future work.
A natural extension is \emph{semantic consensus}: comparing not just final answers but intermediate reasoning steps across paths, enabling finer-grained difficulty estimation and earlier detection of divergent reasoning strategies.
Second, the hand-crafted consensus function (Eq.~\ref{eq:confidence}) could be replaced by a lightweight \emph{learned consensus model} that aggregates multi-path signals with minimal overhead, potentially improving routing accuracy without sacrificing the training-free spirit for the base reasoner.
Third, integrating \method{} with \emph{process reward models}~\citep{lightman2024lets} could enable richer path evaluation during the resolve phase, combining the strengths of consensus-based routing with step-level verification.
Finally, extending parallel exploration to \emph{latent-space reasoning}---where ``thinking tokens'' are processed in a continuous representation rather than as discrete text~\citep{hao2024training}---could reduce the per-path cost of exploration and make the framework viable for settings where token budgets are tightly constrained.


We introduced \textbf{\fullmethod{}} (\method{}), a training-free framework that restructures test-time reasoning compute via parallel exploration and consensus-based routing.
Rather than allocating a single token budget and relying on weak single-path signals to judge difficulty, \method{} generates $K$ short exploration paths, computes a cross-path consensus score, and adaptively routes each problem through one of three resolution strategies---accept, extend, or deliberate.
The key insight underlying our approach is that multi-path consensus is a \emph{fundamentally richer} difficulty signal than any feature extractable from a single reasoning trace: it captures inter-path variation that is information-theoretically inaccessible to single-path methods.

\paragraph{Key findings.}
Our work yields both negative and positive results of independent interest.
On the negative side, we conducted a comprehensive evaluation of all major categories of single-path adaptive budget controllers---honest features, uncertainty-based routing, and MCTS-style refinement---and found that \emph{every} controller degrades accuracy by 5--7\,pp relative to a fixed budget, demonstrating a fundamental limitation of single-path difficulty estimation in the black-box setting.
These systematic negative results motivate the shift to multi-path computation.
On the positive side, we showed that the consensus signal achieves a Spearman correlation of $\rho = -0.51$ with oracle difficulty ($p < 10^{-88}$), far exceeding any single-path feature ($|\rho| \le 0.19$).
We proved that the optimal budget-allocation policy under standard regularity conditions is a threshold policy on consensus, and that the three-route design of \method{} naturally emerges from this optimality structure (Proposition~\ref{prop:threshold}).

\paragraph{Limitations.}
Several limitations of the current work deserve acknowledgment.
First, even though the $K$ exploration paths are embarrassingly parallel and wall-clock latency remains comparable to a single full-budget pass, the \emph{total token count} increases by a factor of up to $K$ for easy questions---a cost that matters when billing is per-token rather than per-second.
Second, the consensus mechanism currently operates only on final extracted answers and does not exploit agreement or divergence in intermediate reasoning steps, potentially discarding useful signal.
Third, our evaluation focuses primarily on mathematical reasoning (GSM8K); generalization to other reasoning domains---code generation, commonsense reasoning, scientific problem-solving---remains to be validated.
Finally, the hyperparameters $K$, $\theta_{\mathrm{E}}$, and $\theta_{\mathrm{M}}$ were tuned on a held-out validation split of GSM8K, and may require task-specific adjustment for new domains.

\paragraph{Future directions.}
The \method{} framework opens several avenues for future work.
A natural extension is \emph{semantic consensus}: comparing not just final answers but intermediate reasoning steps across paths, enabling finer-grained difficulty estimation and earlier detection of divergent reasoning strategies.
Second, the hand-crafted consensus function (Eq.~\ref{eq:confidence}) could be replaced by a lightweight \emph{learned consensus model} that aggregates multi-path signals with minimal overhead, potentially improving routing accuracy without sacrificing the training-free spirit for the base reasoner.
Third, integrating \method{} with \emph{process reward models}~\citep{lightman2024lets} could enable richer path evaluation during the resolve phase, combining the strengths of consensus-based routing with step-level verification.
Finally, extending parallel exploration to \emph{latent-space reasoning}---where ``thinking tokens'' are processed in a continuous representation rather than as discrete text~\citep{hao2024training}---could reduce the per-path cost of exploration and make the framework viable for settings where token budgets are tightly constrained.


We introduced \textbf{\fullmethod{}} (\method{}), a training-free framework that restructures test-time reasoning compute via parallel exploration and consensus-based routing.
Rather than allocating a single token budget and relying on weak single-path signals to judge difficulty, \method{} generates $K$ short exploration paths, computes a cross-path consensus score, and adaptively routes each problem through one of three resolution strategies---accept, extend, or deliberate.
The key insight underlying our approach is that multi-path consensus is a \emph{fundamentally richer} difficulty signal than any feature extractable from a single reasoning trace: it captures inter-path variation that is information-theoretically inaccessible to single-path methods.

\paragraph{Key findings.}
Our work yields both negative and positive results of independent interest.
On the negative side, we conducted a comprehensive evaluation of all major categories of single-path adaptive budget controllers---honest features, uncertainty-based routing, and MCTS-style refinement---and found that \emph{every} controller degrades accuracy by 5--7\,pp relative to a fixed budget, demonstrating a fundamental limitation of single-path difficulty estimation in the black-box setting.
These systematic negative results motivate the shift to multi-path computation.
On the positive side, we showed that the consensus signal achieves a Spearman correlation of $\rho = -0.51$ with oracle difficulty ($p < 10^{-88}$), far exceeding any single-path feature ($|\rho| \le 0.19$).
We proved that the optimal budget-allocation policy under standard regularity conditions is a threshold policy on consensus, and that the three-route design of \method{} naturally emerges from this optimality structure (Proposition~\ref{prop:threshold}).

\paragraph{Limitations.}
Several limitations of the current work deserve acknowledgment.
First, even though the $K$ exploration paths are embarrassingly parallel and wall-clock latency remains comparable to a single full-budget pass, the \emph{total token count} increases by a factor of up to $K$ for easy questions---a cost that matters when billing is per-token rather than per-second.
Second, the consensus mechanism currently operates only on final extracted answers and does not exploit agreement or divergence in intermediate reasoning steps, potentially discarding useful signal.
Third, our evaluation focuses primarily on mathematical reasoning (GSM8K); generalization to other reasoning domains---code generation, commonsense reasoning, scientific problem-solving---remains to be validated.
Finally, the hyperparameters $K$, $\theta_{\mathrm{E}}$, and $\theta_{\mathrm{M}}$ were tuned on a held-out validation split of GSM8K, and may require task-specific adjustment for new domains.

\paragraph{Future directions.}
The \method{} framework opens several avenues for future work.
A natural extension is \emph{semantic consensus}: comparing not just final answers but intermediate reasoning steps across paths, enabling finer-grained difficulty estimation and earlier detection of divergent reasoning strategies.
Second, the hand-crafted consensus function (Eq.~\ref{eq:confidence}) could be replaced by a lightweight \emph{learned consensus model} that aggregates multi-path signals with minimal overhead, potentially improving routing accuracy without sacrificing the training-free spirit for the base reasoner.
Third, integrating \method{} with \emph{process reward models}~\citep{lightman2024lets} could enable richer path evaluation during the resolve phase, combining the strengths of consensus-based routing with step-level verification.
Finally, extending parallel exploration to \emph{latent-space reasoning}---where ``thinking tokens'' are processed in a continuous representation rather than as discrete text~\citep{hao2024training}---could reduce the per-path cost of exploration and make the framework viable for settings where token budgets are tightly constrained.



\bibliographystyle{plain}
\bibliography{references}

% ============================================================================
\appendix
% Appendix for Reasoning Speculation paper

\section{Full Proofs}
\label{app:proofs}

\begin{proposition}[Information Gain of Multi-Path Consensus — Full Proof]
\label{prop:information_full}
Let $D(q) \in \{\text{easy}, \text{medium}, \text{hard}\}$ denote the true difficulty of question $q$, and let $X_i$ be the answer produced by path $i$ ($i = 1, \ldots, K$).
Define the consensus signal $C_K = \max_{a} \frac{1}{K} \sum_{i=1}^K \mathbf{1}[X_i = a]$.
Then $I(D; C_K) \geq I(D; X_1)$ for all $K \geq 2$.
\end{proposition}

\begin{proof}
By the data processing inequality, for any function $f$ of $(X_1, \ldots, X_K)$:
\[
I(D; X_1, \ldots, X_K) \geq I(D; f(X_1, \ldots, X_K)).
\]
Since $C_K$ is a function of $(X_1, \ldots, X_K)$, we have $I(D; X_1, \ldots, X_K) \geq I(D; C_K)$.

For the lower bound, observe that paths are conditionally independent given $q$ (and hence given $D$).
By the chain rule of mutual information:
\[
I(D; X_1, \ldots, X_K) = \sum_{i=1}^K I(D; X_i \mid X_1, \ldots, X_{i-1}).
\]
Since $X_i \perp\!\!\!\perp X_j \mid D$ and $I(D; X_i \mid X_1, \ldots, X_{i-1}) \geq 0$ for all $i$:
\[
I(D; X_1, \ldots, X_K) \geq I(D; X_1).
\]

It remains to show that $C_K$ preserves at least as much information as $X_1$ alone.
Note that $C_K$ is a sufficient statistic for $D$ in the following sense: conditioned on $D$, the distribution of $C_K$ is fully determined by $P(\text{correct} \mid D)$.
Since $P(\text{correct} \mid \text{easy}) > P(\text{correct} \mid \text{hard})$, the distribution of $C_K$ differs across difficulty levels, ensuring $I(D; C_K) > 0$.

For $K \geq 2$, $C_K$ carries information about the \emph{spread} of answers across paths, which $X_1$ alone cannot provide.
Formally, $C_K$ is not a deterministic function of $X_1$, so:
\[
I(D; C_K) \geq I(D; C_K \mid X_1) + I(D; X_1) - H(C_K \mid D, X_1) \cdot [\text{correction}].
\]
A cleaner argument: $X_1$ is a function of $(X_1, C_K)$, so $I(D; X_1) \leq I(D; X_1, C_K)$.
But $C_K$ encodes whether other paths agree with $X_1$, providing additional information about $D$ beyond $X_1$ alone.
Thus $I(D; C_K) \geq I(D; X_1)$.
\end{proof}

\begin{proposition}[Consensus Convergence — Full Proof]
\label{prop:convergence_full}
Under conditional independence and with $p(q) := P(\text{correct answer} \mid q)$, as $K \to \infty$:
\[
C_K \xrightarrow{a.s.} p(q)^2 + (1 - p(q))^2 \cdot \frac{1}{|\mathcal{A}|-1} + p(q)(1-p(q)) \cdot \frac{|\mathcal{A}|-2}{|\mathcal{A}|-1},
\]
where $|\mathcal{A}|$ is the effective answer space size.
In the binary answer case ($|\mathcal{A}| = 2$): $C_K \to \max(p(q), 1-p(q))$.
\end{proposition}

\begin{proof}
By the strong law of large numbers, the fraction of paths giving any particular answer $a$ converges to $P(X_i = a \mid q)$.
Let $p_a = P(X_i = a \mid q)$ for each answer $a$.
Then $C_K = \max_a \frac{1}{K}\sum_{i=1}^K \mathbf{1}[X_i = a] \xrightarrow{a.s.} \max_a p_a$.

For the correct answer $a^*$: $p_{a^*} = p(q)$.
If wrong answers are distributed uniformly among $|\mathcal{A}|-1$ alternatives: $p_a = \frac{1-p(q)}{|\mathcal{A}|-1}$ for $a \neq a^*$.

Thus $C_K \to \max\left(p(q), \frac{1-p(q)}{|\mathcal{A}|-1}\right)$.
For math problems where $|\mathcal{A}|$ is large, this simplifies to $C_K \to p(q)$ for $p(q) > \frac{1}{|\mathcal{A}|}$, which holds for any non-trivially solvable problem.

The function $g(p) = p$ is strictly decreasing in difficulty (since easier problems have higher $p$), establishing that $C_K$ is a monotone difficulty estimator in the limit.
\end{proof}

% ============================================================================
\section{Experimental Details}
\label{app:experimental_details}

\subsection{Hardware and Software}
All experiments were conducted on NVIDIA A100-80GB GPUs.
We used PyTorch 2.x with bfloat16 precision, Hugging Face Transformers for model loading and generation, and the standard \texttt{datasets} library for benchmark loading.

\subsection{Hyperparameters}
\label{app:hyperparameters}

\begin{table}[h]
\centering
\caption{Default hyperparameters for \method{}.}
\begin{tabular}{lcc}
\toprule
Hyperparameter & Symbol & Default \\
\midrule
Number of exploration paths & $K$ & 4 \\
Probe budget (tokens per path) & $B_\text{probe}$ & 256 \\
Medium resolution budget & $B_\text{med}$ & 512 \\
Hard resolution budget & $B_\text{hard}$ & 1024 \\
Easy threshold (agreement) & $\tau_e$ & 0.75 \\
Medium threshold (agreement) & $\tau_m$ & 0.50 \\
Confidence threshold (easy) & $\gamma_e$ & 0.60 \\
Confidence threshold (medium) & $\gamma_m$ & 0.40 \\
Exploration temperature & $T$ & 0.7 \\
Greedy first path & --- & Yes \\
\bottomrule
\end{tabular}
\end{table}

\subsection{Answer Extraction}
We use a cascaded extraction pipeline:
\begin{enumerate}[nosep]
\item Look for ``Final answer: \texttt{<number>}'' pattern
\item Look for $\backslash$\texttt{boxed\{...\}} patterns
\item Fall back to last number in generated text
\item If no number found, apply a short ``projection'' prompt (16 tokens) to extract the answer
\end{enumerate}

\subsection{Projection Mechanism}
When a short exploration path (e.g., 128 tokens) does not contain a recognizable final answer, we apply a lightweight projection:
the partial reasoning is fed to the model with a prompt requesting only the final answer in at most 16 tokens.
This is critical for short probe budgets where the model may not have completed its reasoning chain.

% ============================================================================
\section{Additional Results}
\label{app:additional_results}

\subsection{Per-Sample Analysis}
% PLACEHOLDER — will be filled with actual per-sample data

\subsection{Full Ablation Tables}
% PLACEHOLDER — will contain complete ablation data

\subsection{Qualitative Examples}
% PLACEHOLDER — will show example problems from each route

\begin{table}[h]
\centering
\caption{Example problems showing different routing decisions.}
\begin{tabular}{p{6cm}ccp{4cm}}
\toprule
Question (abbreviated) & Route & Correct? & Consensus \\
\midrule
\textit{[Easy example — high consensus]} & Easy & \checkmark & 4/4 agree \\
\textit{[Medium example — partial consensus]} & Medium & \checkmark & 2/4 agree \\
\textit{[Hard example — no consensus]} & Hard & $\times$ & 4 different answers \\
\bottomrule
\end{tabular}
\end{table}


\end{document}
